\documentclass[11pt,a4paper]{article}

\usepackage[T1]{fontenc}
\usepackage[utf8]{inputenc}
\usepackage[italian]{babel}
\usepackage[margin=2.3cm]{geometry}
\usepackage{lmodern,microtype}
\usepackage{amsmath,amssymb}
\usepackage{booktabs,tabularx,longtable,array,multirow}
\usepackage{enumitem}
\usepackage[dvipsnames]{xcolor}
\usepackage{tikz}
\usetikzlibrary{arrows.meta,positioning,calc,fit,backgrounds}
\usepackage[hyphens]{url}
\usepackage{hyperref}

\definecolor{acc}{RGB}{31,84,140}
\definecolor{sens}{RGB}{214,232,247}
\definecolor{mem}{RGB}{222,240,216}
\definecolor{dec}{RGB}{252,231,207}
\definecolor{val}{RGB}{248,219,219}
\definecolor{mod}{RGB}{233,224,245}
\definecolor{fix}{RGB}{232,232,232}
\hypersetup{colorlinks=true,linkcolor=acc,urlcolor=acc,citecolor=acc,
  pdftitle={Il cervello a sezioni},pdfauthor={Evolutionary nets}}

\setlength{\parindent}{0pt}
\setlength{\parskip}{5pt}
\emergencystretch=2em
\setlist{nosep,leftmargin=1.4em}
\renewcommand{\arraystretch}{1.25}

% Le voci di ogni scheda.
\newcommand{\voce}[1]{\par\smallskip\noindent{\bfseries\color{acc}#1}\par\nopagebreak}
\newcommand{\gene}[1]{\texttt{#1}}
\newcommand{\lesione}[1]{\textsf{\bfseries#1}}
\newcolumntype{Y}{>{\raggedright\arraybackslash}X}
\newenvironment{struttura}{\par\small\noindent\tabularx{\textwidth}{@{}>{\bfseries}L{2.7cm}Y@{}}\toprule}{\bottomrule\endtabularx\par\medskip}
\newcolumntype{L}[1]{>{\raggedright\arraybackslash}p{#1}}

\tikzset{
  sec/.style={draw=black!65,rounded corners=2pt,minimum height=11mm,minimum width=20mm,
              align=center,font=\small,line width=.5pt},
  flow/.style={-{Stealth[length=2mm]},line width=.6pt,black!75},
  soft/.style={flow,dashed},
  teach/.style={-{Stealth[length=2mm]},line width=.9pt,red!70!black},
  lab/.style={font=\scriptsize,fill=white,inner sep=1pt},
}

\title{\textbf{Il cervello a sezioni}\\[2mm]
\large Com'è fatto, a che cosa serve ogni parte, che cosa succederebbe senza\\
e a quali parti del cervello vero corrisponde}
\author{Evolutionary nets}
\date{30 ottobre 2026}

\begin{document}
\maketitle

\begin{abstract}
\noindent Questo documento descrive il cervello a sezioni come è scritto oggi in
\texttt{evonets/brains/pile/}, con le formule e i valori di partenza del codice.
Per ogni sezione dice che cosa fa, perché serve, che cosa succederebbe in teoria
senza di lei e a quale struttura del cervello somiglia. Distingue sempre ciò che
\emph{il codice fa}, ciò che \emph{ci si aspetta in teoria} e ciò che \emph{è stato
misurato}: non sempre coincidono, e la sezione~\ref{sec:stato} dice dove.
\end{abstract}

\tableofcontents
\newpage

% ---------------------------------------------------------------------------
\section{L'idea}

\subsection{Il problema da cui nasce}

Le reti a forma libera del progetto evolvono bene, ma il test di capacità
(\texttt{docs/SPEC\_CAPACITA.md}) ha mostrato che \textbf{non imparano dentro la
vita}, anche con la regola giusta. Le cause misurate sono tre:
\begin{itemize}
\item il cambiamento utile (0,3\% per frutto) affoga nel rumore dell'apprendimento
  hebbiano (10\% in una vita);
\item quando arriva il dolore il cervello non sa più \emph{che cosa} ha mangiato: il
  frutto è fuori vista;
\item il critico sui sensi grezzi peggiora le cose: l'occhio che si adatta durante
  l'avvicinamento sembra una delusione.
\end{itemize}
La conclusione è che non mancano calcoli ma \textbf{parti con ruoli diversi}: una
memoria di ciò che si è appena visto, una descrizione stabile, un posto dove la
scelta viene fatta e imparata.

\subsection{Il principio: si fissa la struttura, non il contenuto}

\begin{center}
\begin{tabularx}{\textwidth}{lYY}
\toprule
& \textbf{Fissato da noi} & \textbf{Emerge dall'esperienza} \\
\midrule
Architettura & sezioni, numero di neuroni, chi parla con chi & \\
Velocità & quanto cambia e quanto dura ogni parte & \\
Regole & una famiglia di regole per sezione & \\
Contenuto & & che cosa rappresenta ogni neurone, quale colore è buono, quale scelta conviene \\
\bottomrule
\end{tabularx}
\end{center}

Nessun neurone nasce con un significato. Tutte le 67 costanti del cervello sono
geni (\texttt{Genes}): nella fase A valgono per tutti i valori di partenza, nella
fase B le tara l'evoluzione, insieme ai pesi di partenza e alle connessioni di C.
L'evoluzione \emph{tara} le regole, non le inventa: è la lezione del progetto Cyber
Rodent (Elfwing, Uchibe, Doya, Christensen, 2011).

\subsection{Il compito}

Il cervello vive nell'\textbf{arena delle pile}, un recinto di $500\times500$ con 6
esseri e 8 pile. Le pile sono rosse o gialle: una ricarica (+25 di energia, piacere),
l'altra toglie 8 e fa male. \textbf{Ogni 300 secondi i ruoli si scambiano} senza che
niente lo segnali. Ogni 200 secondi cambia l'angolo in cui compaiono le pile. Si
mangia solo mordendo, e mordere costa. A energia zero si muore: nella fase A si
rinasce con il cervello di partenza, senza ciò che si era imparato. Il mondo fa un
passo ogni $\Delta t = 0{,}05$ s.

Il compito è fatto apposta per chiedere ciò che le reti di prima non sanno fare:
imparare in fretta \emph{quale} colore conviene e ricredersi quando il mondo cambia.

% ---------------------------------------------------------------------------
\section{Lo schema}\label{sec:schema}

Le sezioni sono dodici più due meccanismi trasversali (l'arbitraggio e la soglia
sulla sorpresa). Il percorso completo dai sensi ai motori, con i cicli, è nella
sezione~\ref{sec:flusso}. La figura~\ref{fig:attivita} mostra chi manda \emph{attività} a chi
mentre l'essere agisce; la figura~\ref{fig:apprendimento} mostra chi \emph{insegna} a
chi. Sono due schemi separati perché rispondono a due domande diverse: il primo a
``come decide?'', il secondo a ``come impara?''.

\begin{figure}[ht]
\centering
\begin{tikzpicture}[x=1cm,y=1cm]
\node[sec,fill=fix,minimum width=14mm] (W) at (0.6,6) {Sensi\\31};
\node[sec,fill=sens] (A) at (3.0,6) {\textbf{A} interfaccia\\52 canali};
\node[sec,fill=sens] (S) at (3.0,9) {\textbf{S} attenzione\\mappa di 9 posizioni};
\node[sec,fill=sens] (B) at (6.0,6) {\textbf{B} nucleo\\16 prototipi};
\node[sec,fill=mem] (D) at (5.6,3.6) {\textbf{D} mem. di lavoro\\10 canali, 4 s};
\node[sec,fill=mem] (C) at (9.6,3.6) {\textbf{C} mem. veloce\\256 neuroni, 3 min};
\node[sec,fill=mem] (L) at (9.6,1.2) {\textbf{L} mem. lunga\\stessa forma di C};
\node[sec,fill=val] (O) at (9.2,9) {\textbf{O} valore\\degli oggetti};
\node[sec,fill=val] (K) at (12.2,9) {\textbf{K} contesti\\4 cassetti di O};
\node[sec,fill=dec] (E) at (12.2,6) {\textbf{E} selezione\\4 programmi + morso};
\node[sec,fill=fix,minimum width=17mm] (F) at (15.2,6) {\textbf{F} motore\\4 neuroni};
\node[sec,fill=fix,minimum width=17mm] (M) at (15.2,1.2) {Corpo\\e mondo};

\draw[flow] (W) -- (A);
\draw[flow] (S) -- node[lab,right] {sceglie la fovea} (A);
\draw[flow] (A) -- (B);
\draw[flow] (B) -- (E);
\draw[flow] (E) -- (F);
\draw[flow] (F) -- (M);
% la fovea verso D e C
\draw[flow] ([xshift=-5mm]A.south) |- node[lab,pos=.88,above] {fovea} (D.west);
\draw[flow] ([xshift=5mm]A.south) -- ++(0,-3.3) -| node[lab,pos=.3,above] {che cos'è l'oggetto} ([xshift=-7mm]C.south);
% B verso C
\fill (8.9,6) circle (.6mm);
\draw[flow] (8.9,6) -- ([xshift=-7mm]C.north);
% D verso C, E, O
\draw[flow] (D) -- (C);
\fill (7.7,3.6) circle (.6mm);
\draw[flow] (7.7,3.6) -- (7.7,4.9) -- (11.4,4.9) -- ([xshift=-8mm]E.south);
\draw[flow] (7.7,4.9) -- (7.7,7.6) -- (8.7,7.6) -- ([xshift=-5mm]O.south);
\node[lab] at (7.7,7.1) {ciò che ho preso};
% C e L verso E
\draw[flow] (C.east) -| (12.2,5.45);
\draw[soft] (C) -- node[lab,right] {solo a riposo} (L);
\draw[flow] (L.east) -| ([xshift=8mm]E.south);
% O, K, S
\draw[flow] (K) -- (O);
\draw[flow] (O) -- node[lab,above=1pt,align=center] {valore, importanza,\\incertezza} (S);
\draw[flow] ([xshift=5mm]O.south) -- (9.7,7.4) -- (11.4,7.4) -- ([xshift=-8mm]E.north);
% lascialo
\draw[flow] (13.9,6.55) -- (13.9,10.3) -- node[lab,above] {``lascialo'': l'oggetto viene ignorato per 3 s} (3.0,10.3) -- (S.north);
% il morso non passa da F
\draw[flow] (13.6,5.45) -- node[lab,pos=.45] {morso} (13.6,1.2) -- (M.west);
% il mondo risponde
\draw[flow] (M.south) -- ++(0,-.5) -| (W.south);
\end{tikzpicture}
\caption{Il flusso dell'attività. Da sinistra a destra: i sensi diventano canali (A),
la situazione (B), i ricordi (D, C, L) e una scelta (E), che F trasforma in movimento.
O colora gli oggetti di valore; S decide quale guardare. Azzurro: percezione. Verde:
memorie. Rosso: valore. Arancio: decisione. Grigio: parti fisse.}
\label{fig:attivita}
\end{figure}

\begin{figure}[ht]
\centering
\begin{tikzpicture}[x=1cm,y=1cm]
\node[sec,fill=fix] (M) at (0.9,4) {Corpo\\piacere, dolore,\\energia spesa};
\node[sec,fill=val] (R) at (3.9,4) {$r$ valore\\del momento};
\node[sec,fill=val] (T) at (7.5,4) {$\hat\delta$ soglia e scala\\(la sorpresa che insegna)};
\node[sec,fill=mod] (U) at (7.5,1.2) {\textbf{U} incertezza\\``adesso'' / ``il solito''};
\node[sec,fill=mem,minimum width=33mm] (C) at (12.9,5.7) {\textbf{C} uscite verso E};
\node[sec,fill=dec,minimum width=33mm] (E) at (12.9,4.2) {\textbf{E} pesi dell'abitudine};
\node[sec,fill=mem,minimum width=33mm] (D) at (12.9,2.7) {\textbf{D} cancello};
\node[sec,fill=dec,minimum width=33mm] (AR) at (12.9,1.2) {Arbitraggio: pesi di\\E+L, C, O$\to$S; accelera C};
\node[sec,fill=val] (OK) at (3.9,7) {\textbf{O} e \textbf{K}\\Rescorla-Wagner};
\node[sec,fill=sens] (B) at (3.9,1.2) {\textbf{B} prototipi};
\node[sec,fill=val,dashed] (V) at (8.2,7) {\textbf{V} tre critici\\imparano ma non insegnano};
\draw[teach] (M) -- (R);
\draw[teach] (R) -- node[lab,above] {$\delta$} (T);
\draw[teach] (T.east) -- (10.3,4) |- (C.west);
\draw[teach] (T.east) -- (10.3,4) |- (E.west);
\draw[teach] (T.east) -- (10.3,4) |- (D.west);
\draw[teach] (T) -- (U);
\draw[flow] (U.east) -- (AR.west);
\draw[teach] (M.north) |- node[lab,pos=.72,above] {esito di ogni episodio} (OK.west);
\draw[teach] (R.south) -- node[lab,right] {$|\delta|$: quanto conta il momento} (B.north);
\draw[soft] (R.north east) ++(-3mm,0) -- (V.south west);
\draw[soft] (V.south) ++(-.7,0) -- node[lab,right] {\gene{critic\_mix} = 0} (7.5,4.55);
\end{tikzpicture}
\caption{Il flusso dell'apprendimento. Il corpo è l'unico maestro. In rosso i segnali
che insegnano davvero oggi. O e K imparano direttamente dagli esiti; C, E e il cancello
di D imparano dalla sorpresa filtrata $\hat\delta$; U sposta il peso delle decisioni. Il
tratteggio indica il critico V, che con i valori attuali non influisce su niente
(sezione~\ref{sec:V}).}
\label{fig:apprendimento}
\end{figure}

\subsection{L'ordine dei calcoli in un passo}

L'ordine conta: ogni passaggio usa ciò che i precedenti hanno appena aggiornato.
\begin{enumerate}
\item Gli episodi di piacere o dolore appena finiti insegnano a K e a O.
\item S sceglie l'oggetto da guardare; se è cambiato, E deciderà subito.
\item A prepara i canali, compresa la fovea sull'oggetto scelto.
\item Il corpo dà il valore del momento $r$.
\item B, D e C si aggiornano; O legge i due valori ``ciò che guardo'' e ``ciò che ho preso''.
\item I critici V imparano; da $r$ si ricava $\delta$, poi la soglia dà $\hat\delta$ e U si aggiorna.
\item C impara e dimentica; a riposo L si avvicina a C.
\item E somma i voti, tenendo conto dell'arbitraggio, impara e, se è il momento, sceglie.
\item F traduce il programma in sterzo e spinta (è calcolato dall'arena, fuori dal cervello).
\end{enumerate}

\subsection{Le scale di tempo}

Il motivo più profondo per avere tante sezioni è che \textbf{una sola velocità non può
bastare}. Ciò che impara in fretta dimentica in fretta; ciò che è stabile non si adatta.
Ogni sezione occupa una scala diversa (McClelland, McNaughton, O'Reilly, 1995).

\begin{center}
\begin{tabularx}{\textwidth}{llY}
\toprule
\textbf{Scala} & \textbf{Dove} & \textbf{Come dura} \\
\midrule
0,05 s -- 0,5 s & A, B, decisioni di E e S & attività che segue l'ingresso \\
2 s & tracce di E, C, cancello & ``chi ha appena agito'' resta marcato \\
4 s & D & attività che resta accesa, nessuna sinapsi cambia \\
20 s / 200 s & U, soglia & medie della sorpresa: ``adesso'' e ``il solito'' \\
3 minuti & C & sinapsi che imparano in un colpo e sbiadiscono \\
minuti -- vita & O, K & sinapsi che cambiano a ogni episodio \\
vita & B, L, E & sinapsi lente \\
generazioni & geni & costanti, pesi di partenza, connessioni di C \\
\bottomrule
\end{tabularx}
\end{center}

% ---------------------------------------------------------------------------
\section{L'anatomia completa}\label{sec:anatomia}

Questa sezione è la mappa da tenere a fianco delle schede. I numeri sono quelli del codice e
si riferiscono a un essere solo; tutte le costanti sono i valori di partenza dei geni.

\subsection{Che cosa vuol dire ``neurone'' qui}

Non ci sono impulsi: ogni neurone è un numero che cambia nel tempo. Non tutti però sono
fatti allo stesso modo. Nel cervello a sezioni esistono \textbf{sei tipi di unità}.

\begin{center}
\small
\begin{tabularx}{\textwidth}{@{}L{3.1cm}YL{2.7cm}@{}}
\toprule
\textbf{Tipo} & \textbf{Come calcola} & \textbf{Dove} \\
\midrule
Canale & copia un valore, non calcola & A, fovea \\
Filtro & segue l'ingresso con ritardo: $x\leftarrow\kappa x+(1-\kappa)\,\text{ingresso}$ & A, D, B (uscita) \\
Prototipo & risponde tanto più quanto l'ingresso è \emph{vicino} a un punto salvato nei suoi pesi: misura una distanza, non fa una somma pesata & B \\
A soglia & somma pesata $z$, poi resta acceso solo chi vince una competizione & C \\
Lineare & somma pesata $z=\sum w_i x_i$, senza soglia & V, O, E (punteggi), F \\
Sigmoide & $\sigma(z)=1/(1+e^{-z})$: un valore tra 0 e 1 & cancello di D, morso di E \\
\bottomrule
\end{tabularx}
\end{center}

Ogni sinapsi che impara ha quasi sempre accanto una \textbf{traccia di eleggibilità}: un
numero che dice ``questa sinapsi ha appena lavorato''. Quando arriva la sorpresa, cambia
solo la sinapsi con la traccia alta. È il modo con cui un segnale globale, che arriva
uguale a tutti, modifica soltanto le connessioni giuste.

\subsection{Quanto è grande ogni sezione}

\begin{center}
\small
\begin{tabularx}{\textwidth}{@{}lL{3.2cm}L{3.0cm}rY@{}}
\toprule
& \textbf{Unità} & \textbf{Connessioni fisse} & \textbf{Pesi che imparano} & \textbf{Altri numeri tenuti} \\
\midrule
A & 52 canali, di cui 11 con filtro & --- & 0 & 11 medie lente \\
S & 9 posizioni & 7 pesi uguali per tutte & 0 & oggetto attento, novità di 14 oggetti \\
B & 16 prototipi & --- & 1664 & attività, medie delle vittorie \\
D & 10 memorie, 1 cancello & 10 copie dalla fovea & 10 & 10 tracce \\
C & 256 neuroni & 1536 dagli ingressi & 1280 & 1280 tracce \\
L & nessuna: solo sinapsi & --- & 1280 & --- \\
O & 2 neuroni di uscita & --- & 56 & 56 incertezze, 6 tracce, episodio \\
K & 4 contesti, 1 attivo & --- & 0 & prove, usi, conteggi \\
E & 5 uscite & --- & 145 & 145 tracce \\
V & 3 critici & --- & 888 & 888 tracce \\
U & 1 unità & --- & 0 & 3 medie \\
F & 4 neuroni motori & 56 pesi (13 non nulli) & 0 & deriva, verso di fuga \\
\midrule
\textbf{Totale} & & & \textbf{5323} & \\
\bottomrule
\end{tabularx}
\end{center}

I 5323 pesi sono ciò che l'essere impara in vita e perde quando muore. Il genoma ha 9201
numeri: 67 costanti, 31 valori innati legati ai sensi, 56 pesi di F, 832 prototipi di partenza,
7168 interruttori per le connessioni possibili di C ($256\times28$; una c'è se il suo gene
supera 0,5), 145 pesi di partenza di E, 888 di V, 14 di O.

\subsection{Tutte le connessioni}

Ogni riga è un fascio di fibre. \emph{Fissa} significa che non cambia in vita (tutte le
fisse sono però geni). \emph{Sorpresa} significa che impara con la sorpresa filtrata $\hat\delta$.

\begin{center}
\small
\begin{tabularx}{\textwidth}{@{}L{2.5cm}rL{2.6cm}Y@{}}
\toprule
\textbf{Da $\to$ a} & \textbf{Numero} & \textbf{Regola} & \textbf{Note} \\
\midrule
Sensi, fovea $\to$ A & 41 & fissa, copia & più 11 canali ``è appena salito'' \\
A $\to$ B & $16\times52$ & competitiva & sono i prototipi: si confrontano con l'ingresso \\
B $\to$ previsione di A & $52\times16$ & errore locale & serve a B per sapere che cosa segue \\
Fovea $\to$ D & 10 & fissa, una a una & passa solo a cancello aperto \\
Fovea $\to$ cancello & 10 + 1 & sorpresa; fissa & 10 imparano; ``nuovo'' ha peso 8 \\
Identità in fovea $\to$ C & \multirow{3}{*}{1536} & fissa, peso 1 & canali 1--6 \\
Identità in D $\to$ C & & fissa, peso 0,5 & \\
B $\to$ C & & fissa, peso 0,1 & il contesto \\
C $\to$ E (voto veloce) & $5\times256$ & sorpresa & decade in 3 minuti, limite $\pm0{,}3$ \\
C $\to$ E (voto di L) & $5\times256$ & media di C & cambia solo a riposo \\
B, D, O, 1 $\to$ E & $5\times29$ & sorpresa & 16 + 10 + 2 + 1 ingressi; limite $\pm5$ \\
Fovea, D, C, energia, B, O, 1 $\to$ V & $3\times296$ & TD & 10 + 10 + 256 + 1 + 16 + 2 + 1 ingressi \\
Identità $\to$ O & $2\times7\times4$ & Rescorla-Wagner & piacere e dolore, 6 + sfondo, 4 cassetti \\
O $\to$ S & 9 posizioni & fissa & tre pesi: 2, 0,5, 0,5 \\
S, E $\to$ F & $4\times14$ & fissa & il programma combinato con la posizione \\
Raggi $\to$ F & 3 & fissa & il riflesso dei muri \\
\bottomrule
\end{tabularx}
\end{center}

% ---------------------------------------------------------------------------
\section{Sezione per sezione}

Ogni scheda ha la stessa forma. \emph{Struttura} dice di che cosa è fatta la sezione:
neuroni, connessioni e stato, con il nome che hanno nel codice. \emph{Come funziona}
elenca che cosa fa a ogni passo di tempo, nell'ordine esatto. \emph{Perché
serve} spiega il problema che risolve. \emph{Se mancasse} è il ragionamento teorico;
dove esiste una lesione nel codice è indicata la sua lettera. \emph{Nel cervello} dà
l'analogia più vicina: è un'ispirazione, non un'equivalenza. \emph{Misure} riporta i
dati disponibili con la loro origine.

% --- A ---
\subsection{A --- Interfaccia sensoriale}

\voce{Struttura}
\begin{struttura}
Unità & 52 canali, senza pesi. 31 grezzi + 11 ``è appena salito'' + 10 di fovea. \\
I 31 grezzi & 9 occhi (rosso, giallo, blu $\times$ sinistra, centro, destra); 12 occhi per i muri (4 colori
  $\times$ 3 direzioni); 3 raggi di distanza; 3 urti; energia; ``mi guardano''; gusto; dolore. Tutti tra 0 e 1. \\
Gli 11 veloci & uno per ognuno dei 9 occhi, per il gusto e per il dolore. I muri e i raggi non li hanno. \\
Stato & \texttt{slow}: 11 medie, una per canale veloce. \\
Geni & \gene{tau\_a} = 0,3 s. \\
Connessioni & riceve i sensi dal mondo e la fovea da S; manda tutti i 52 canali soltanto a B. \\
\end{struttura}

\voce{Come funziona, a ogni passo}
\begin{enumerate}
\item Per ognuno degli 11 canali calcola quanto il segnale supera la sua media recente, e lo
  taglia tra 0 e 1. Una discesa dà 0: conta solo ciò che sale.
\item Aggiorna la media:
\end{enumerate}
\[
\text{veloce}_t=\bigl[x_t-\text{lento}_t\bigr]_0^1,
\qquad
\text{lento}_{t+1}=\kappa\,\text{lento}_t+(1-\kappa)\,x_t,
\qquad \kappa=e^{-\Delta t/0{,}3\,\text{s}}.
\]
\begin{enumerate}[resume]
\item Unisce sensi, veloci e fovea in un vettore di 52 numeri, in quest'ordine.
\end{enumerate}
Quindi una luce che si accende dà un impulso nel canale veloce, che sparisce in circa 0,3 s
anche se la luce resta. La fovea ha una scheda sua (sezione~\ref{sec:fovea}).

\voce{Perché serve}
Separa ``com'è'' da ``che cosa è cambiato''. Un cervello che legge solo i valori
assoluti non distingue l'inizio di un evento dal suo perdurare, e gli eventi sono
proprio ciò che va imparato. Inoltre isola il resto del cervello dai sensori: per
cambiare sensori basta cambiare A.

\voce{Se mancasse}
Il cervello non avrebbe più nessun ingresso, quindi non è una lesione possibile. Si
può ragionare su pezzi di A. Senza i canali veloci B vedrebbe solo stati e avrebbe più
difficoltà a segnare gli istanti in cui succede qualcosa. Senza la fovea non ci sarebbe
il concetto di ``l'oggetto'': i sensi direbbero che c'è del rosso a sinistra, non che
\emph{quella cosa} è rossa, grande e vicina, e D, C e O non avrebbero niente da leggere.

\voce{Nel cervello}
Retina e talamo sensoriale. I canali lenti e veloci ricordano le vie sostenute e
transitorie della visione (parvocellulare e magnocellulare).

% --- S ---
\subsection{S --- Attenzione}

\voce{Struttura}
\begin{struttura}
Unità & 9 neuroni-posizione. Coprono i $120^\circ$ davanti al muso, $13{,}3^\circ$ ciascuno. Un
  neurone rappresenta una \emph{direzione}, non un oggetto: gli oggetti vengono e vanno, i neuroni no. \\
Ingressi di ogni neurone & dell'oggetto più vicino in quella direzione (gli altri sono nascosti): vicinanza,
  novità, se è già l'oggetto attento, se è appena stato lasciato, e i tre numeri di O. \\
Connessioni & 7 pesi fissi, uguali per tutti i 9 neuroni. Nessuna sinapsi impara. \\
Stato & \texttt{target}: l'oggetto attento. \texttt{inhib}, \texttt{inhib\_t}: l'oggetto lasciato e quanto manca.
  \texttt{novelty}: quanto è nuovo ognuno dei 14 posti oggetto. \texttt{prev\_vis}: chi si vedeva prima.
  \texttt{new}: il canale ``nuovo'' della fovea. \texttt{changed}: l'oggetto è cambiato. \\
Geni & \gene{sal\_near} 2; \gene{sal\_novelty} 1,5; \gene{sal\_stick} 1; \gene{sal\_inhibit} 4; \gene{sal\_value} 2;
  \gene{sal\_import} 0,5; \gene{sal\_unc} 0,5; \gene{sal\_temperature} 0,4; \gene{novelty\_tau} e \gene{new\_tau} 0,5 s;
  \gene{inhibit\_time} 3 s; \gene{decide\_every} 0,5 s. \\
\end{struttura}

\voce{Come funziona, a ogni passo}
\begin{enumerate}
\item La novità di ogni oggetto cala con $\tau=0{,}5$ s; chi è appena comparso riceve +1.
\item Il tempo di inibizione scende. Se l'oggetto attento non si vede più, \texttt{target} diventa
  $-1$. Il canale \texttt{new} cala con $\tau=0{,}5$ s.
\item Si costruisce la mappa: ogni oggetto in vista va nella posizione del suo angolo. Se in una
  posizione ce ne sono più di uno, resta il più vicino.
\item Ogni posizione occupata calcola il suo punteggio:
\end{enumerate}
\begin{align*}
z={}&\underbrace{2\,\text{vicino}+1{,}5\,\text{novità}+1\,[\text{è l'attento}]-4\,[\text{appena lasciato}]}_{\text{innato}}\\
&+\underbrace{\ell\,(2\,v_O+0{,}5\,i_O)+0{,}5\,u_O}_{\text{da O}},
\qquad \ell=1-0{,}7\,k,
\end{align*}
dove $v_O$, $i_O$, $u_O$ sono valore, importanza e incertezza dell'oggetto e $k$ viene da U.
\begin{enumerate}[resume]
\item I punteggi diventano probabilità: $P_i\propto e^{z_i/0{,}4}$.
\item Se è il momento di decidere (ogni 0,5 s) o non c'è un oggetto attento, si estrae una
  posizione a caso con quelle probabilità. Con temperatura 0,4, una differenza di punteggio pari a 1
  vale circa 12 volte in probabilità. Se l'oggetto è diverso da prima, \texttt{new} = 1.
\item Se l'oggetto è cambiato si alza \texttt{changed}, ed E decide subito.
\end{enumerate}
Oltre a O e a U, su S agisce soltanto E, con il programma ``lascialo'', che mette l'oggetto in
\texttt{inhib} per 3 s.

\voce{Perché serve}
Risolve il \textbf{problema del merito}. Se un essere guarda tre pile insieme e poi
sente dolore, quale ha causato il dolore? Con un oggetto alla volta la domanda ha una
risposta. Inoltre lega la caratteristica all'oggetto: ``rosso'' e ``vicino'' riguardano
la stessa cosa. Infine porta il sapere nell'azione: vedere prima ciò che conviene è un
modo di scegliere che non passa da E.

Il punteggio ``importanza'' fa sì che anche il pericolo attiri lo sguardo; quello
``incertezza'' spinge a guardare le cose di cui non si sa ancora abbastanza, che è il
modo per imparare che cosa valgono.

\voce{Se mancasse (lesione \lesione{S}: resta solo la salienza innata)}
L'essere guarderebbe ciò che è vicino e nuovo, buono o cattivo che sia. O continuerebbe
a imparare, ma il suo sapere arriverebbe alle azioni solo attraverso i pesi di E, che
sono lenti. Ci si aspetta che la scelta cada molto, e che però il crollo dopo
un'inversione sia più piccolo: non avendo imparato a preferire un colore, c'è meno da
disimparare.

Togliendo l'intera mappa di attenzione, senza un oggetto attento, si perderebbe anche
la fovea, e con lei il funzionamento di D, C e O.

\voce{Nel cervello}
Il collicolo superiore: una mappa di posizioni, dove vince un punto e la testa o gli
occhi si orientano. Il punteggio combinato è il concetto di \emph{mappa di priorità}
della corteccia parietale e frontale (Fecteau e Munoz, 2006; Bisley e Goldberg, 2010).
``Appena lasciato'' è l'\emph{inibizione del ritorno} (Posner). L'attrazione verso ciò
che è incerto viene da Pearce e Hall (1980).

\voce{Misure}
Con 15 minuti di mondo (\texttt{SPEC\_ARENA\_PILE.md}, 29 ottobre): con O e S che lo
legge 58,2\% di pile buone; con S che non lo legge 52,5\%, quanto senza O. In una misura
precedente, senza S la quota di buone prima dell'inversione era 55\% contro 65\%, con un
crollo minore dopo (47\% contro 37\%): il comportamento previsto.

% --- fovea ---
\subsection{La fovea}\label{sec:fovea}

\voce{Che cos'è}
La fovea \textbf{non è un sensore}: è un elenco di 10 numeri che descrive \emph{un solo
oggetto}, quello scelto da S. Il nome è una metafora. Nel cervello a sezioni ``guardare
una cosa'' non vuol dire voltarsi: vuol dire riempire questi 10 numeri con quella cosa.
Se S non ha scelto niente, i primi 9 valgono zero.

\voce{Com'è fatta}
\begin{center}
\small
\begin{tabularx}{\textwidth}{@{}clYl@{}}
\toprule
\textbf{N.} & \textbf{Canale} & \textbf{Come si calcola} & \textbf{Che cosa dice} \\
\midrule
1--3 & rosso, giallo, blu & 0 o 1 & che cos'è \\
4--5 & grande, piccolo & 0 o 1 (solo le pile) & che cos'è \\
6 & si muove & 1 se un essere va a più del 5\% della velocità massima & che cos'è \\
7 & vicino & $1-\text{distanza}/350$ & dove sta \\
8 & a sinistra & $\text{angolo}/60^\circ$ se è a sinistra, altrimenti 0 & dove sta \\
9 & a destra & $\text{angolo}/60^\circ$ se è a destra, altrimenti 0 & dove sta \\
10 & nuovo & 1 quando l'oggetto cambia, poi cala con $\tau=0{,}5$ s & da quando lo guardo \\
\bottomrule
\end{tabularx}
\end{center}
I primi 6 canali sono l'\textbf{identità} dell'oggetto, gli ultimi 4 il suo
\textbf{rapporto con me}. Una pila ha sempre rosso o giallo, e grande o piccolo. Gli
esseri hanno il blu (la luce sul muso) e, se si muovono, ``si muove''. Quindi il blu
significa di fatto ``è un compagno''. I muri non sono oggetti: non entrano mai nella fovea.

\voce{Un esempio}
Una pila rossa grande, a 140 di distanza e $20^\circ$ a destra, appena scelta da S:
\[
\text{fovea}=(\underbrace{1,\ 0,\ 0}_{\text{rossa}},\ \underbrace{1,\ 0}_{\text{grande}},\
\underbrace{0}_{\text{ferma}},\ \underbrace{0{,}6}_{\text{vicino}},\
\underbrace{0}_{\text{sinistra}},\ \underbrace{0{,}33}_{\text{destra}},\
\underbrace{1}_{\text{nuovo}}).
\]

\voce{Da dove vengono i numeri}
Questo è il punto più importante, e il più diverso da un cervello vero. \textbf{La fovea
non viene calcolata dagli occhi.} Il mondo passa direttamente al cervello un elenco degli
oggetti in vista (funzione \texttt{objects\_of}), fino a 13 (8 pile e 5 compagni), con le
loro caratteristiche esatte. Un oggetto è in vista se è entro $60^\circ$ dal muso e a
meno di 350; un compagno anche se è rivolto verso di me. I 9 canali degli occhi sono
un'altra via, molto più grossolana: ogni occhio dice soltanto quanto è vicino il colore
più vicino in un settore di $40^\circ$.

Vedere, separare gli oggetti dallo sfondo e riconoscerli è quindi \textbf{un lavoro già
fatto} prima di arrivare al cervello. È una scelta precisa: il progetto studia che cosa
si impara a fare degli oggetti, non come si impara a vederli.

\voce{Dove va}
\begin{itemize}
\item In \textbf{A}, come ultimi 10 canali, e da lì in B.
\item In \textbf{D}: ne resta una copia quando l'oggetto è nuovo.
\item In \textbf{C}: solo i 6 canali dell'identità, perché il codice non cambi con la
  distanza e lo stesso oggetto accenda gli stessi neuroni a ogni incontro.
\item In \textbf{O}: l'identità diventa la traccia a cui andrà il merito di un esito.
\item In \textbf{V}: tutti i 10 canali.
\item In \textbf{F} no, formalmente: F riceve ``dove sta l'oggetto'' direttamente da S.
  È la stessa informazione dei canali 7--9, presa dalla stessa fonte.
\end{itemize}
E non legge la fovea direttamente: la conosce attraverso B, D e i due neuroni di O.

\voce{Perché serve}
La fovea trasforma una scena in \emph{una cosa alla volta}, e con questo risolve due
problemi insieme. Il primo è il merito: dopo un esito si sa a chi attribuirlo. Il secondo
è il legame: i canali degli occhi dicono che c'è del rosso e c'è del grande, ma non che
la stessa pila è rossa e grande. Nella fovea queste proprietà stanno per forza insieme,
perché appartengono all'unico oggetto presente. Separare identità e posizione permette
inoltre di imparare ``che cos'è'' senza confonderlo con ``dove si trova''.

\voce{Se mancasse}
Con i soli occhi, il cervello saprebbe che c'è del rosso a sinistra, ma due pile dello
stesso settore sarebbero una cosa sola e la grandezza non si vedrebbe affatto, perché
nessun occhio la misura. D non avrebbe niente da tenere, C niente da riconoscere e O
niente a cui legare gli esiti. Anche F perderebbe un bersaglio preciso da inseguire.
Nell'arena con le pile buone per grandezza il compito diventerebbe impossibile.

\voce{Nel cervello}
La fovea vera è la piccola zona centrale della retina, con la massima risoluzione, e per
usarla bisogna muovere gli occhi. \textbf{Qui è diverso in due modi.} S può scegliere un
oggetto ovunque nel ventaglio di $120^\circ$ senza muoversi, e questo somiglia di più
all'\emph{attenzione nascosta}, che si sposta senza spostare lo sguardo (Posner, 1980).
Inoltre la descrizione è perfetta, mentre quella vera è il risultato di decine di
millisecondi di calcolo nella corteccia visiva. Il concetto più vicino è il \emph{file
dell'oggetto} (Kahneman, Treisman, Gibbs, 1992): una scheda mentale che tiene insieme le
proprietà di un oggetto mentre lo si segue. Lo stesso nome ``fovea'' resta qui per
comodità.

% --- B ---
\subsection{B --- Nucleo stabile}

\voce{Struttura}
\begin{struttura}
Unità & 16 neuroni-prototipo. Ogni neurone è un punto nello spazio dei 52 canali di A. \\
Pesi & \texttt{proto}: $16\times52=832$. Sono il punto salvato da ogni neurone. Partono a caso tra 0 e 0,3
  (dal genoma), quindi non significano niente. \\
 & \texttt{pred}: $52\times16=832$. Da B prevedono il prossimo ingresso di A. Partono da zero. \\
Ingressi & tutti i 52 canali di A, e $|\delta|$ come velocità di apprendimento. Nessun ritorno dall'alto. \\
Uscite & 16 valori di attività (\texttt{state\_b}) verso C (peso 0,1), E e V (peso 0,2). \\
Stato & \texttt{state\_b}, \texttt{prev\_b}; \texttt{win\_rate}: quanto spesso ha vinto ogni neurone di recente. \\
Geni & \gene{b\_winners} = 3 (struttura); \gene{conscience} 2; \gene{eta\_b} 0,05 al s; \gene{eta\_predict} 0,5 al s;
  \gene{predict} 0,5; \gene{surprise\_boost} 2; \gene{tau\_b} 0,2 s; \gene{win\_window} 30 s. \\
\end{struttura}

\voce{Come funziona, a ogni passo}
\begin{enumerate}
\item \textbf{Distanza.} Ogni neurone misura quanto l'ingresso $a$ dista dal suo punto, e riceve un aiuto
  se ha vinto meno della sua quota $1/16$ (la ``coscienza''):
\[
d_j=\lVert \text{proto}_j-a\rVert^2-2\,\bigl(\tfrac{1}{16}-\text{vittorie}_j\bigr).
\]
\item \textbf{Competizione.} Si ordinano le distanze. Vincono i 3 più vicini; la distanza del quarto,
  $d_{(4)}$, fa da metro:
\[
y_j=\Bigl[1-\frac{d_j}{d_{(4)}}\Bigr]_0^1,\quad\text{poi}\ y\leftarrow y/\max y.
\]
  Il più vicino vale sempre 1. Tutti gli altri, dal quarto in poi, valgono 0.
\item \textbf{Errore di previsione.} Si confronta l'ingresso con ciò che \texttt{pred} prevedeva a
  partire da \texttt{prev\_b}, cioè dall'attività di due passi prima: $\text{err}=a-\text{pred}\cdot\text{prev\_b}$.
\item \textbf{La previsione impara}, con la regola delta: $\Delta\text{pred}=0{,}5\,\Delta t\;\text{err}\otimes\text{prev\_b}$.
\item \textbf{I vincitori imparano}: si spostano verso l'ingresso e un po' verso ciò che era imprevisto,
  tanto più se il momento è sorprendente ($|\delta|$ viene tagliato a 1):
\[
\Delta\text{proto}_j=0{,}05\,\Delta t\,(1+2|\delta|)\;y_j\,\bigl[(a-\text{proto}_j)+0{,}5\,\text{err}\bigr].
\]
\item \textbf{Vittorie.} \texttt{win\_rate} segue, con $\tau=30$ s, $\frac13$ per chi ha vinto e 0 per gli altri.
\item \textbf{Uscita.} L'attività non salta: segue $y$ con $\tau=0{,}2$ s, cioè
  $\text{state}\leftarrow\text{state}+\frac{\Delta t}{0{,}2}(y-\text{state})$.
\end{enumerate}
Misurando durante una corsa: in media 4,4 neuroni hanno attività sopra 0,01, perché ai 3 vincitori
si sommano quelli che si stanno spegnendo, e il massimo è in media 0,91.

\voce{Perché serve}
I 52 canali cambiano continuamente, ma le situazioni che contano sono poche. B
riassume l'ingresso in un codice stabile e compatto, senza che nessuno gli dica che
cosa rappresentare. Il codice è lo stesso per lungo tempo, e questo è ciò che permette
alle parti che scelgono di imparare su qualcosa di fermo. La previsione lo spinge a
preferire le distinzioni che servono a capire che cosa succederà; il fattore
$1+2|\delta|$ gli fa dedicare più risorse ai momenti che contano.

\voce{Se mancasse (lesione \lesione{B}: E, V e C leggono A)}
E dovrebbe imparare su 52 canali grezzi invece che su 16 stati: più pesi, più lento, più
sensibile al rumore. Inoltre ogni piccolo cambiamento dei sensi cambierebbe l'ingresso
delle sezioni successive, e le generalizzazioni ``questa è la stessa situazione di
prima'' sarebbero più difficili. Per un compito semplice il danno può essere piccolo:
il colore è già scritto chiaramente nella fovea. Il vantaggio di B dovrebbe vedersi nei
compiti in cui la situazione è una combinazione di molti sensi.

\voce{Nel cervello}
La neocorteccia sensoriale e associativa, che organizza da sola le sue rappresentazioni.
Corrisponde all'idea di Doya (1999) che la corteccia impari senza maestro. La
competizione ricorda l'inibizione laterale, la previsione la codifica predittiva (Rao e
Ballard, 1999).

\voce{Misure}
Le lesioni di B, alle ultime misure disponibili, non hanno cambiato il comportamento in
modo chiaro (54--57\% contro 56\%). In una misura precedente B portava poca informazione
sul colore: i prototipi si specializzano su ciò che è frequente (muri, energia), non su
ciò che conta. La scheda va letta con questo limite: oggi il vantaggio di B è più teorico
che dimostrato.

% --- D ---
\subsection{D --- Memoria di lavoro}

\voce{Struttura}
\begin{struttura}
Unità & 10 unità di memoria (\texttt{mem}), una per canale della fovea, e 1 neurone-cancello. \\
Connessioni fisse & ogni canale della fovea va alla sua unità: una a una, peso 1. \\
Pesi & \texttt{w\_gate}: 10, dai canali della fovea al cancello. Partono da zero, limite $\pm10$. \\
Il cancello & è un neurone sigmoide con 12 connessioni: i 10 canali (pesi che imparano), ``nuovo'' una seconda
  volta (peso fisso 8) e una corrente di fondo ($-6$). È uno solo: apre o chiude tutte le 10 unità insieme. \\
Stato & \texttt{mem}: 10 valori. \texttt{e\_gate}: 10 tracce. \\
Geni & \gene{tau\_d} 4 s; \gene{gate\_event} 8; \gene{gate\_bias} $-6$; \gene{eta\_gate} 0,003; \gene{tag} 0,5;
  \gene{trace\_actor} 2 s. \\
Uscite & tutti i 10 valori a E e V; i primi 6, l'identità, a C (peso 0,5) e a O. \\
\end{struttura}

\voce{Come funziona, a ogni passo}
\begin{enumerate}
\item Il cancello impara con la sorpresa del passo prima: $\Delta w_g=0{,}003\,\hat\delta\;e_g$.
\item Calcola la sua apertura:
\[
z=8\,\text{nuovo}+\sum_i w_{g,i}\,\text{fovea}_i-6,\qquad g=\sigma(z).
\]
  A riposo $g=0{,}25\%$; con ``nuovo'' a 1, $g=88\%$. Il cancello supera il 50\% quando ``nuovo''
  supera 0,75: dato che ``nuovo'' cala con $\tau=0{,}5$ s, resta aperto circa 0,15 s, cioè 3 passi.
\item Aggiorna le tracce, solo dai canali della fovea sopra 0,5:
  $e_g\leftarrow e^{-\Delta t/2\,\text{s}}\,e_g+g(1-g)\,[\text{fovea}]_{>0{,}5}$.
\item Aggiorna la memoria, una miscela tra tenere e scrivere:
\[
m_{t+1}=(1-g)\,e^{-\Delta t/4\,\text{s}}\,m_t+g\,\text{fovea}_t.
\]
\end{enumerate}
Il termine $g(1-g)$ è la derivata del cancello: la traccia è alta quando il cancello era indeciso, e
lì un cambiamento dei pesi conta di più. Tenere non usa nessuna sinapsi: il ricordo è attività che cala
del 63\% in 4 s, mentre imparano solo i 10 pesi che decidono \emph{quando} scrivere.

\voce{Perché serve}
Risolve il \textbf{ritardo tra l'azione e l'esito}. Quando si morde, la pila è davanti
al muso un attimo prima, poi l'energia arriva dopo: il gusto o il dolore si sentono
quando la pila può non essere più in vista. Per capire \emph{che cosa} ha fatto bene o
male serve ricordare che cosa si aveva davanti. È esattamente il difetto trovato nel
test di capacità. D risponde alla domanda ``che cosa ho appena preso?'', e da lì nasce
il secondo neurone di O.

La correzione trovata dai test spiega bene il suo funzionamento: con il cancello
socchiuso (5\% a passo) il ricordo spariva in meno di un secondo, perché ogni passo
scriveva sopra un pezzo di fovea vuota. Tenere una cosa significa \emph{non} scrivere.

\voce{Se mancasse (lesione \lesione{D})}
O vedrebbe solo ciò che è nella fovea durante gli ultimi $\sim$0,5 s (la sua traccia):
gli esiti che arrivano più tardi resterebbero senza colpevole e finirebbero sullo
sfondo. C perderebbe l'oggetto ricordato. E non saprebbe di avere appena preso qualcosa.
In compiti con ritardi lunghi il danno sarebbe grave. Nell'arena delle pile lo è meno,
perché la pila presa di solito è ancora in vista, e questo spiega perché la lesione
misuri poco.

\voce{Nel cervello}
La corteccia prefrontale dorsolaterale, che tiene l'informazione con attività persistente
(Goldman-Rakic, 1995). Il cancello imparato viene dal modello di O'Reilly e Frank (2006),
dove i gangli della base decidono quando la prefrontale si aggiorna.

\voce{Misure}
Nessun effetto chiaro alle ultime lesioni (54--57\% contro 56\%).

% --- C ---
\subsection{C --- Memoria veloce}

\voce{Struttura}
\begin{struttura}
Unità & 256 neuroni a soglia. Non hanno pesi propri che imparano in ingresso. \\
Ingressi possibili & 28: i 6 di identità in fovea (peso 1), i 6 di identità in D (peso 0,5), i 16 di B (peso 0,1). \\
Connessioni in ingresso & ogni neurone ne ha esattamente 6, scelte a caso alla nascita tra le 28, tutte di valore 1
  prima del peso dell'ingresso: 1536 in tutto. Non cambiano in vita (\texttt{c\_in}). \\
Inibizione & un'unica inibizione globale, che non è un neurone del codice ma un calcolo: lascia acceso il 5\% dei neuroni. \\
Pesi in uscita & \texttt{u\_c}: $5\times256=1280$, verso le 5 uscite di E. Partono da zero, limite $\pm0{,}3$. \\
Stato & \texttt{code}: 256 valori 0 o 1. \texttt{e\_c}: 1280 tracce. \\
Geni & \gene{c\_sparse} 0,05; \gene{c\_threshold} 0,4; \gene{context} 0,1; \gene{memory} 0,5; \gene{eta\_c} 0,1;
  \gene{c\_max} 0,3; \gene{forget} 180 s; \gene{surprise\_rate} 1; \gene{stick} 1,5. \\
\end{struttura}

\voce{Come funziona, a ogni passo}
\begin{enumerate}
\item \textbf{Ingresso.} Si prepara $h=(\text{identità fovea},\ 0{,}5\,\text{identità D},\ 0{,}1\,\text{B})$.
\item \textbf{Somma.} Ogni neurone somma i suoi 6 ingressi: $z_c=\sum_{j} c_{cj}\,h_j$.
\item \textbf{Competizione.} Con $k=\text{round}(0{,}05\times256)=13$, si trova il valore del tredicesimo
  neurone più eccitato: questa è la soglia. Si accendono quelli sopra la soglia. I pari merito alla
  soglia si accendono solo se, in totale, gli accesi non supererebbero 26. Chi sta sotto 0,4 non si accende mai.
\item \textbf{Apprendimento e oblio}, a ogni passo:
\[
\Delta u=\frac{0{,}1\,(1+U)\,\hat\delta}{1+\#\text{accesi}}\;e-u\,\frac{\Delta t}{180\ \text{s}},
\qquad |u|\le0{,}3 .
\]
\item La traccia cala con $\tau=2$ s.
\item \textbf{Voto.} Il voto di C per ognuna delle 5 uscite è $\sum_c u_{c}\,\text{code}_c$.
\item \textbf{Inerzia.} Finché la traccia dura, i neuroni accesi votano anche con le tracce stesse:
  $1{,}5\,(e\cdot\text{code})/\#\text{accesi}$. Tornando sullo stesso oggetto, riaffiora la scelta.
\item \textbf{Marcatura.} Quando E sceglie, la traccia dei neuroni accesi riceve la deviazione della scelta:
  $e\leftarrow e+(\text{scelta}-\text{attesa})\otimes\text{code}$.
\end{enumerate}

\voce{Il codice, misurato}
Una pila ha due caratteristiche accese (colore e grandezza). I neuroni che le leggono entrambe hanno
$z=2$, quelli che ne leggono una $z=1$, gli altri meno. Misurando su 400 cervelli nati a caso, senza B e con
D vuota, ``rossa e grande'' accende in media 9,9 neuroni, da 2 a 21 secondo il caso. Le sovrapposizioni dicono
che il codice è la combinazione, non la somma delle parti:

\begin{center}
\small
\begin{tabular}{lr}
\toprule
\textbf{Rossa grande confrontata con} & \textbf{Neuroni in comune} \\
\midrule
Rossa piccola (stesso colore) & 14,9\% \\
Gialla grande (stessa grandezza) & 16,1\% \\
Gialla piccola (niente in comune) & 1,9\% \\
Un compagno in movimento & 1,9\% \\
\bottomrule
\end{tabular}
\end{center}

Con la stessa pila anche in D si sale a 12,8 neuroni, il 77,5\% dei quali è lo stesso gruppo. Durante
una corsa, con B attivo, i neuroni accesi mentre si guarda sono in media 13,5: i piccoli contributi di B
rompono i pari merito. Lo stesso calcolo mostra il prezzo del contesto: a parità di oggetto il codice
non è identico.

\voce{Perché serve}
Due cose. La prima è \textbf{riconoscere le combinazioni}. Molti neuroni più degli
ingressi e un'attivazione rada trasformano ``rossa e grande'' in un codice proprio,
diverso da ``rossa e piccola'' e da ``gialla e grande''. Con un solo ingresso per
accendersi, i codici si sovrapponevano seguendo le singole caratteristiche; con la
soglia si accendono solo i neuroni che ricevono più caratteristiche insieme.

La seconda è \textbf{imparare in un colpo}. E è lenta perché deve mediare su molte
esperienze; C impara quasi subito e satura presto, così conta l'ultima esperienza, non
la media di una vita. Il dimenticare in 3 minuti non è un difetto: in un mondo che si
inverte ogni 5 minuti, un ricordo che non sbiadisce diventa un errore. Inoltre C tiene
la marcatura delle scelte per qualche secondo, e finché dura vota per quella scelta
(``inerzia legata all'oggetto''): tornando su un oggetto guardato poco fa riaffiora ciò
che si era deciso per lui. È la memoria di lavoro nelle sinapsi di Mongillo, Barak e
Tsodyks (2008).

\voce{Se mancasse (lesione \lesione{C})}
Sparirebbe la decisione rapida: dopo un'inversione resterebbe solo E, che è lenta, e O
attraverso l'attenzione. Il riadattamento dovrebbe peggiorare. Sparirebbe anche
l'arbitraggio, perché non ci sarebbe più niente a cui spostare il peso. \textbf{Nel codice
la lesione spegne anche L}: L legge lo stesso codice di C, quindi senza C il suo voto è
zero. E nemmeno l'inerzia sull'oggetto resta. Una lesione di C misura quindi più di C.

\voce{Nel cervello}
L'analogia più stretta non è con i mammiferi: è il \textbf{corpo a fungo degli insetti}.
Le cellule di Kenyon ricevono proiezioni casuali e fisse dall'olfatto (Caron e colleghi,
2013), sono molte più degli ingressi, sono tenute rade da un neurone inibitore gigante
(APL) e imparano solo in uscita, con la dopamina. Nei mammiferi il concetto corrisponde
alla separazione dei pattern del giro dentato e, per il ricordo rapido, all'ippocampo
(Marr, 1971).

\voce{Misure}
Una lesione precedente: senza C, 61\% di buone prima dell'inversione contro 65\%. Il C
attuale, con la regola delle cellule di Kenyon, ha portato il massimo di energia: 40,5
al minuto (\texttt{TRAGUARDI\_PILE.md}), ma con un riadattamento peggiore: una memoria
più nitida è più dura da cambiare.

% --- O ---
\subsection{O --- Valore degli oggetti}

\voce{Struttura}
\begin{struttura}
Unità & 2 neuroni di uscita (\texttt{o\_out}): ``quanto vale ciò che guardo'' e ``quanto vale ciò che ho preso''. \\
Pesi & 7 ingressi: le 6 caratteristiche più uno \emph{sfondo}, sempre acceso. Per ogni ingresso un peso per il piacere
  ($w^+$, \texttt{o\_pl}) e uno per il dolore ($w^-$, \texttt{o\_pa}). In tutto 14 pesi per cassetto, 4 cassetti (K),
  quindi 56. Limite $\pm5$. Partono dal genoma (zero nella fase A). \\
Incertezza & per ogni ingresso due numeri (\texttt{o\_assoc}): la sorpresa di adesso e quella solita. Partono da 0,5. \\
Traccia & \texttt{o\_e}: 6 valori, il massimo tra identità in fovea e in D, con coda di 0,5 s. \\
Episodio & \texttt{ep\_on}: se è in corso. \texttt{ep\_e}: la traccia fotografata all'inizio. \texttt{ep\_pl}, \texttt{ep\_pa}: somme
  di piacere e dolore. \\
Geni & \gene{eta\_o} 1,5; \gene{need} 0,8; \gene{pain\_weight} 2; \gene{trace\_o} 0,5 s; \gene{assoc\_rate} 0,3;
  \gene{assoc\_slow} 0,03; \gene{observe} 1. \\
Connessioni & riceve identità da fovea e D, esiti dal corpo, energia dai sensi, il cassetto da K. Manda i 3 numeri
  a S, per ogni oggetto; i 2 neuroni a E e V. \\
\end{struttura}

\voce{Come funziona: a ogni passo legge}
Per un oggetto con caratteristiche $f$ (0 o 1), dal cassetto attivo:
\begin{gather*}
\text{fame}=1+0{,}8\,\bigl(2(1-\text{energia})-1\bigr),\\
v=\sum_f f\,(w^+_f\,\text{fame}-2\,w^-_f),
\qquad
i=\sum_f f\,\bigl([w^+_f]_++[2w^-_f]_+\bigr),
\qquad
u=\sum_f f\,\text{sorpresa}_f .
\end{gather*}
La fame va da 0,2 (carichi) a 1,8 (scarichi). Lo sfondo non entra: è ciò che capita senza oggetto.

\voce{Come funziona: a ogni episodio impara}
\begin{enumerate}
\item Un episodio comincia al primo passo con piacere o dolore. Si fotografa la traccia $e$: sei
  caratteristiche più lo sfondo a 1. È ciò che si aveva in mente \emph{prima} dell'esito.
\item Durante l'episodio si sommano piacere e dolore. Una pila buona dà in tutto 1 in 2 s; una
  cattiva dà $8/25=0{,}32$ di dolore in un passo.
\item Quando il corpo smette di sentire, l'episodio finisce. K sceglie il cassetto.
\item Si calcolano gli errori, uno per piacere e uno per dolore: $\text{err}^+=\text{piacere}-\sum_f w^+_f e_f$.
\item Si calcola la velocità. La volatilità è il rapporto tra sorpresa di adesso e solita, limitata tra 0,5 e 4:
\[
\alpha_f=\min\bigl(1{,}5\cdot\text{volatilità}_f,\ 1{,}9\bigr).
\]
\item Si aggiornano i pesi, dividendo per quanta traccia c'è perché la correzione totale non superi l'errore:
\[
\Delta w^{\pm}_f=\alpha_f\,\frac{e_f}{1+\lVert e\rVert^2}\;\text{err}^{\pm}.
\]
\item Si aggiorna l'incertezza con la grandezza $s=\min(|\text{err}^+|+|\text{err}^-|,2)$: la sorpresa di adesso
  con peso 0,3 e quella solita con peso 0,03, in proporzione alla traccia.
\end{enumerate}
Se il maestro è attivo, ogni allievo ripete i passi 4--7 con la traccia e l'esito del maestro.

\voce{Perché serve}
È la sezione che ha reso possibile l'apprendimento del colore. Rispetto a E e C ha tre
vantaggi.

\textbf{Le caratteristiche competono.} L'errore è unico per tutte: se ``rosso'' prevede
già l'esito, ``grande'' non guadagna niente. È il \emph{blocking}, e significa che il
merito va a ciò che prevede davvero. Lo stesso meccanismo trova da solo se nel mondo
conta il colore o la grandezza.

\textbf{Il valore non dipende dal comportamento.} Un critico classico impara quanto vale
una situazione \emph{con il comportamento di adesso}. O impara che cosa fa una cosa
quando la si prende, e questo non cambia con il modo di muoversi.

\textbf{La fame si applica leggendo.} Da sazi il cibo vale meno senza dover imparare
niente: la conoscenza e il bisogno sono separati. Il dolore pesa il doppio del piacere:
perdere conta più che guadagnare.

\voce{Se mancasse (lesione \lesione{O})}
S tornerebbe alla sola salienza innata, quindi non ci sarebbe più preferenza dello
sguardo, e il sapere sui colori sarebbe solo nei voti di C ed E, imparati da una
sorpresa grezza e meno mirata. Ci si aspetta un crollo della scelta verso il caso. Si
perderebbero anche la ricerca di ciò che non si conosce e la differenza tra sazi e
affamati nel giudicare le cose.

\voce{Nel cervello}
La corteccia orbitofrontale e l'amigdala, che tengono il valore atteso degli oggetti e
lo aggiornano con gli esiti (Schoenbaum e colleghi, 2009). Il fatto che sazietà cambi
il valore senza apprendimento è la \emph{svalutazione} studiata in orbitofrontale.
Separare piacere e dolore rispecchia i circuiti distinti di ricompensa e punizione.
La regola è Rescorla e Wagner (1972); l'incertezza Pearce e Hall (1980); l'accelerazione
quando il mondo è volatile Behrens e colleghi (2007). Il maggior peso del dolore è
l'avversione alla perdita (Kahneman e Tversky, 1979).

\voce{Misure}
È l'effetto di lesione più chiaro: 58,2\% con O, 52,5\% senza. O preferisce la pila buona
nell'84\% degli esseri dopo 2 minuti e nel 90\% dopo 5; 30 secondi dopo un'inversione
solo nel 13\%; dopo 5 minuti torna all'89\%.

% --- K ---
\subsection{K --- Contesti}

\voce{Struttura}
\begin{struttura}
Unità & 4 neuroni-contesto, di cui uno solo attivo (\texttt{k\_on}). \\
Pesi & nessuno suo. Ogni contesto possiede un cassetto completo di O: è il contesto che seleziona quale gruppo di 14 pesi
  si legge e si scrive. \\
Stato & \texttt{k\_ev}: la prova di ogni contesto, cioè il suo errore recente. \texttt{k\_used}: quali sono in uso (all'inizio
  solo il primo). \texttt{k\_usual}: l'errore solito. \texttt{k\_since}: episodi dall'ultimo cambio.
  \texttt{k\_stale}: da quanti episodi ogni contesto non è attivo. \texttt{o\_prior}: i valori innati. \\
Geni & \gene{k\_rate} 0,3; \gene{k\_stick} 0,15; \gene{k\_new} 1,5; \gene{k\_wait} 4 episodi. \\
\end{struttura}

\voce{Come funziona, a ogni episodio finito}
\begin{enumerate}
\item Ogni contesto in uso dice quale errore avrebbe fatto: $E_k=|\text{piacere}-\sum w^+_k e|+|\text{dolore}-\sum w^-_k e|$.
\item Aggiorna la sua prova, media con peso 0,3: $\text{prova}_k\leftarrow\text{prova}_k+0{,}3\,(\min(E_k,2)-\text{prova}_k)$.
\item \textbf{Cambio.} Se il contesto con la prova migliore batte quello attivo di almeno 0,15, diventa lui l'attivo.
\item \textbf{Reclutamento.} Sia $\text{limite}=1{,}5\cdot\max(\text{errore solito},0{,}05)$. Si recluta un contesto nuovo solo se
  tutte e tre sono vere: la prova dell'attivo supera il limite; nessun contesto spiega l'ultimo esito, cioè tutti gli $E_k$
  superano il limite; sono passati almeno 4 episodi dall'ultimo cambio.
\item Il posto è il primo libero o, se sono tutti usati, quello fermo da più tempo. Il suo cassetto riparte dai valori innati.
\item \textbf{Solo ora} impara O, e soltanto nel cassetto scelto.
\item L'errore solito si aggiorna con peso 0,03.
\end{enumerate}

\voce{Perché serve}
Il mondo ha uno \textbf{stato nascosto}: quale colore è buono adesso. Non si vede, si
deduce dagli esiti. Senza K, O riscrive gli stessi pesi a ogni inversione: il sapere
vecchio viene distrutto per fare posto al nuovo, e quando il mondo torna com'era va
reimparato da capo. Con K il cassetto vecchio resta intatto e ci si torna in pochi
episodi. È una risposta alla \emph{interferenza catastrofica}: i ricordi vengono
separati per contesto invece che sommati.

\voce{Se mancasse (lesione \lesione{K}: un contesto solo)}
Al primo incontro con un regime nuovo non cambierebbe nulla, perché nessun cassetto può
già saperlo. La differenza si vedrebbe solo \emph{al ritorno} di un regime già vissuto:
ci si impiegherebbe ogni volta quanto la prima. Si perderebbe anche il modo di rispondere
alla domanda ``il mondo è cambiato o ho sbagliato io?''.

\voce{Nel cervello}
L'orbitofrontale come mappa di stati nascosti del compito (Wilson e colleghi, 2014) e
l'ippocampo che lega i ricordi al contesto. Il reclutamento di una causa nuova è la teoria
delle cause latenti (Gershman, Blei, Niv, 2010): spiega perché una paura estinta ritorna
se si torna nel contesto originale. Il fatto che prima si decida il contesto e poi si
impari è anche l'idea del modello COIN (Heald, Lengyel, Wolpert, 2021).

\voce{Misure}
Sul banco di prova non si vede, perché pochi vivono abbastanza da rivedere un regime. Tra
chi sopravvive: nel primo minuto dopo il ritorno di un regime, 54--57\% con K contro
49--51\% senza; il crollo quasi sparisce. Al primo cambio, 41\% in entrambi i casi, come
previsto. Un limite: il 27\% degli esseri riempie i 4 cassetti per rumore.

% --- L ---
\subsection{L --- Memoria lunga}

\voce{Struttura}
\begin{struttura}
Unità & nessuna. L è un secondo fascio di sinapsi dagli stessi 256 neuroni di C alle stesse 5 uscite di E. \\
Pesi & \texttt{u\_l}: 1280. Partono da zero; non hanno limite proprio, perché inseguono \texttt{u\_c}. \\
Geni & \gene{eta\_l} 0,01 al s; \gene{rest\_energy} 0,6. \\
\end{struttura}

\voce{Come funziona, a ogni passo}
\begin{enumerate}
\item Controlla se è a riposo: energia sopra 0,6 \emph{e} gusto sopra 0,1, cioè sta consumando una pila.
\item Se sì, ogni peso si avvicina al corrispondente peso di C. Se no, non cambia niente:
\[
\Delta u_L=0{,}01\,\Delta t\,\mathbb{1}_{\text{riposo}}\,(u_C-u_L).
\]
\item Il voto $\sum u_L\,\text{code}$ si somma a quello di E nell'abitudine.
\end{enumerate}
Non usa tracce e non usa la sorpresa: copia. Servono 100 secondi di riposo perché un peso percorra il 63\% della
distanza da C, e il riposo è raro: per questo in una vita breve L cambia pochissimo.

\voce{Perché serve}
Tiene ciò che vale \emph{sempre}. Poiché è lenta, ciò che si inverte di continuo si
annulla e ciò che resta vero si accumula: un filtro che separa le regolarità stabili
dalle circostanze. Così C può essere velocissima e dimenticare, senza che tutto vada
perso. È l'idea dei sistemi di apprendimento complementari. In questo mondo il colore
si inverte di continuo; ciò che resta vero è per esempio che i compagni non si mangiano.

\voce{Se mancasse (lesione \lesione{L})}
Dopo 3 minuti ciò che non è stato rinforzato sparirebbe del tutto: ogni regolarità
andrebbe reimparata. Nel mondo attuale il danno teorico è piccolo, perché c'è poco di
stabile tra gli esiti e il tempo di riposo è poco, e con $\eta_l$ così basso L accumula
pochissimo in una vita di qualche minuto. La lesione dovrebbe pesare di più in vite
lunghe.

\voce{Nel cervello}
La neocorteccia, che secondo McClelland e colleghi (1995) assorbe lentamente i ricordi
dall'ippocampo, soprattutto a riposo e nel sonno. \textbf{La differenza è grande}: il
cervello vero riproduce i ricordi alternati ai vecchi, L invece copia soltanto; la
riproduzione vera non è ancora stata fatta.

\voce{Misure}
Nessun effetto chiaro.

% --- E ---
\subsection{E --- Selezione dell'azione}

\voce{Struttura}
\begin{struttura}
Unità & 5 uscite: 4 per i programmi (avvicinati, allontanati, lascialo, esplora) e 1 per il morso. \\
Ingressi diretti & 29: 16 di B, 10 di D, 2 di O, 1 costante. Non legge la fovea, né C direttamente in questo fascio. \\
Pesi & \texttt{w\_e}: $5\times29=145$. Partono dal genoma (zero nella fase A), limite $\pm5$. \\
Altri ingressi & i voti di C (\texttt{u\_c}), di L (\texttt{u\_l}) e l'inerzia, tutti da 256 neuroni. \\
Stato & \texttt{e\_e}: 145 tracce. \texttt{program}, \texttt{bite}: la scelta in corso. \texttt{pi}, \texttt{p\_bite}: le probabilità.
  \texttt{clock}: il tempo dall'ultima decisione. \\
Geni & \gene{eta\_e} 0,01; \gene{temperature} 0,5; \gene{stick} 1,5; \gene{tag} 0,5; \gene{trace\_actor} 2 s;
  \gene{decide\_every} 0,5 s; \gene{arb\_fast} 2; \gene{arb\_slow} 0,5. \\
\end{struttura}

\voce{Come funziona, a ogni passo}
\begin{enumerate}
\item \textbf{Somma i voti.} Per ognuna delle 5 uscite, con $f$ i 29 ingressi e $k$ da U:
\[
z=(1-0{,}5\,k)\,\underbrace{(w\cdot f+u_L\cdot\text{code})}_{\text{abitudine: E + L}}
+(1+2\,k)\,\underbrace{u_C\cdot\text{code}}_{\text{memoria veloce}}
+\text{inerzia di C}.
\]
\item \textbf{Impara}, anche se non decide: $\Delta w=0{,}01\,\hat\delta\,e/(1+\lVert f\rVert^2)$, e la traccia cala con $\tau=2$ s.
\item Se non è il momento di decidere, il programma resta quello di prima: fine.
\item \textbf{Sceglie il programma.} Se non sta guardando niente, il programma in corso riceve +1,5. Poi softmax con $T=0{,}5$:
  $\pi_i\propto e^{z_i/0{,}5}$, e si estrae a caso.
\item \textbf{Sceglie il morso}, a parte: $P=\sigma(z_{\text{morso}}/0{,}5)$. All'inizio tutti i pesi sono zero, quindi morde al 50\%.
\item \textbf{Marca.} Per ogni uscita la deviazione è scelta meno attesa, e va sulle tracce:
  $e\leftarrow e+(\text{scelta}-\pi)\,f_{\text{marcati}}$. Di B e D marcano solo i neuroni sopra 0,5; O e la costante sempre.
\item Se il programma è ``lascialo'', S abbandona l'oggetto.
\end{enumerate}
Il programma non nomina un colore: ``avvicinati'' vuol dire ``verso l'oggetto attento''. Per questo
lo stesso programma vale per qualsiasi cosa S abbia scelto. ``Lascialo'' ed ``esplora'' producono lo stesso movimento: si
distinguono perché il primo agisce su S.

\voce{Perché serve}
Qualcuno deve fare una scelta sola. I voti di C, L, D e O possono essere in disaccordo, e
l'azione del corpo deve essere una. E è il punto dove tutto converge. La sua lentezza è
una proprietà: costruisce abitudini sicure, non si fa spostare da un singolo episodio.
La traccia usa \emph{la scelta rispetto all'attesa} perché questo è ciò che rende la
regola un gradiente della ricompensa attesa (Williams, 1992): scegliere una cosa già
quasi certa insegna poco, sorprendersi di una scelta rara insegna molto. Per questo
l'esplorazione non è un dettaglio: senza il caso nella scelta, la regola non ha niente
da cui imparare. Il morso è separato perché è un'altra decisione: guardare una cosa e
mangiarla non sono lo stesso atto.

Le \emph{tracce strette} evitano il rumore che bloccava le reti di prima: una sinapsi si
marca solo se il neurone da cui parte è acceso sopra 0,5.

\voce{Se mancasse}
Non c'è una lesione nel codice, perché senza E non esiste comportamento: F riceverebbe
sempre lo stesso programma. Togliendo solo le sinapsi che imparano resterebbero i voti di
C, ma non più un modo lento e stabile di scegliere: ci si riadatterebbe in fretta ma
senza mai consolidare niente. Togliendo solo il caso ($T\to0$) si sceglierebbe sempre la
stessa cosa e l'apprendimento si fermerebbe.

\voce{Nel cervello}
I gangli della base, e in particolare lo striato: ricevono da tutta la corteccia, le
popolazioni di uscita competono e ne passa una. L'apprendimento è modulato dalla
dopamina, come in Doya (1999). La regola a tre fattori (attività prima, attività dopo,
dopamina) è quella delle sinapsi cortico-striatali. La temperatura è l'esplorazione che
Doya (2002) lega alla noradrenalina.

\voce{Misure}
L'attore-critico classico, che ha E e V ma non le altre sezioni, sceglie al 50,4\%, cioè
a caso, contro il 57,8\% attuale. I pesi da O a E restano però quasi a zero: il sapere di
O arriva alle azioni soprattutto attraverso S, non attraverso E.

% --- V ---
\subsection{V --- Valore, il critico}\label{sec:V}

\voce{Struttura}
\begin{struttura}
Unità & 3 neuroni lineari, uno per orizzonte: 2, 8 e 30 s. Tutti e tre leggono gli stessi ingressi. \\
Ingressi & 296: fovea (10), D (10), C (256), energia (1), B con peso 0,2 (16), O (2), costante (1). \\
Pesi & \texttt{w\_v}: $3\times296=888$. Partono dal genoma (zero nella fase A). \\
Stato & \texttt{e\_v}: 888 tracce. \texttt{v\_prev}: le 3 previsioni del passo prima. \texttt{fresh}: appena nato. \\
Geni & \gene{h\_short} 2, \gene{h\_mid} 8, \gene{h\_long} 30 s; \gene{eta\_v} 0,5; \gene{eta\_vc} 5 per gli ingressi da C;
  \gene{trace\_critic} 1 s; \gene{critic\_mix} 0; \gene{critic\_b} 0,2. \\
Uscite & verso il cervello, nessuna con i valori attuali (vedi sotto). La previsione media si vede solo sullo schermo. \\
\end{struttura}

\voce{Come funziona, a ogni passo}
\begin{enumerate}
\item Per ogni critico, con $\gamma=e^{-\Delta t/h}$ calcola la previsione $V=w\cdot f$ e l'errore TD:
\[
\delta_V=r+\gamma\,V-V_{\text{prima}},\qquad |\delta_V|\le3 .
\]
\item Aggiorna i pesi, normalizzando per non divergere. Gli ingressi da C imparano 10 volte più in fretta:
\[
\Delta w=\frac{\eta\,\delta_V}{1+\lVert f\rVert^2+e\cdot f}\;e .
\]
\item Aggiorna la traccia: $e\leftarrow\gamma\,e^{-\Delta t/1\,\text{s}}\,e+f$.
\item Salva la nuova previsione e fa la media dei tre errori, $\bar\delta_V$.
\end{enumerate}

\voce{Il punto importante: oggi V non insegna}
La sorpresa usata dal resto del cervello è
\[
\delta=r+\text{\gene{critic\_mix}}\cdot(\bar\delta_V-r),
\]
e \textbf{il valore di partenza di \gene{critic\_mix} è 0}: quindi $\delta=r$, il valore
grezzo del momento. V impara, ma il suo errore non arriva a nessuno. È stato verificato:
con e senza la lesione \lesione{V}, a parità di seme, posizioni e pesi di E restano
identici per tutti i 1500 passi provati. Il gene c'è perché la scelta sia dell'evoluzione.

Il motivo è misurato: dando peso al critico si peggiora (a peso 1, 2,30 pile al minuto
contro 3,54). Quando un oggetto entra nella fovea la previsione scende invece di salire,
perché quei momenti arrivano spesso subito dopo aver finito di mangiare: il critico
prevede il valore della situazione \emph{con questo comportamento}, non dell'oggetto.

\voce{Perché servirebbe}
L'idea è trasformare il premio in una \emph{previsione}. Senza critico si impara solo
quando il premio arriva. Con il critico l'indizio che lo precede diventa un premio a sua
volta: si può imparare da lontano, e un premio già previsto smette di insegnare. Tre
orizzonti servono perché ``quanto vale'' dipende da quanto lontano si guarda: mordere una
cattiva è male subito, restare in una zona buona conviene a lungo.

\voce{Se mancasse (lesione \lesione{V})}
\emph{Oggi}: niente, per quanto visto sopra. \emph{In teoria}, se V insegnasse: si
tornerebbe a un cervello che impara solo sull'istante del premio e che continua a
``sorprendersi'' di premi ovvi. Anche i compiti con lunghe attese diventerebbero
difficili.

\voce{Nel cervello}
Lo striato ventrale e i neuroni a dopamina dell'area tegmentale ventrale e della
sostanza nera, che codificano l'errore di previsione (Schultz, Dayan, Montague, 1997).
Orizzonti diversi corrispondono a zone diverse dello striato (Tanaka e colleghi, 2004).

% --- U ---
\subsection{U --- Incertezza e arbitraggio}

\voce{Struttura}
\begin{struttura}
Unità & 1 sola: $U$, un numero tra 0 e 3. Da lei $k=U/3$, tra 0 e 1. \\
Pesi & nessuno. Non impara. \\
Stato & \texttt{d\_fast}: $|\hat\delta|$ medio su 20 s. \texttt{t\_slow}: $|\hat\delta|$ medio su 200 s. Dalla soglia, \texttt{d\_slow}: $|\delta|$ medio su 200 s. \\
Geni & \gene{u\_fast} 20 s; \gene{u\_slow} 200 s; \gene{surprise\_rate} 1; \gene{arb\_fast} 2; \gene{arb\_slow} 0,5; \gene{att\_shift} 0,7. \\
Uscite & $k$ verso E (pesi dei voti) e S; $U$ verso C (velocità). \\
\end{struttura}

\voce{Come funziona, a ogni passo}
\begin{enumerate}
\item Aggiorna le due medie della sorpresa che insegna.
\item Ne fa il rapporto:
\[
U=\Bigl[\frac{\overline{|\hat\delta|}_{20\,\text{s}}}{\overline{|\hat\delta|}_{200\,\text{s}}}-1\Bigr]_0^3 .
\]
\item Il passo dopo, $k=U/3$ cambia i pesi delle decisioni. Con $k=1$, cioè il massimo:
\end{enumerate}
\begin{center}
\small
\begin{tabular}{lcc}
\toprule
& \textbf{Con $k=0$} & \textbf{Con $k=1$} \\
\midrule
Voto di C & $\times1$ & $\times3$ \\
Voti di E e L & $\times1$ & $\times0{,}5$ \\
Valore e importanza da O in S & $\times1$ & $\times0{,}3$ \\
Velocità di apprendimento di C ($U=3$) & $\times1$ & $\times4$ \\
\bottomrule
\end{tabular}
\end{center}
In una corsa di 2 minuti la media è $U\approx1{,}2$, non zero: le sorprese sopra la soglia sono frequenti
anche quando il mondo non cambia.

\voce{Perché serve}
Risponde a ``quanto posso fidarmi di ciò che so?''. Se le sorprese sono le solite, il
mondo è come sempre e conviene l'abitudine; se improvvisamente aumentano, le abitudini
sono scadute. Allora il peso passa a chi impara in fretta. È un cambio di \emph{chi
decide}, non un aumento del caso: nelle versioni precedenti U aggiungeva esplorazione ed
era solo rumore.

U conta solo le sorprese sopra la soglia: con tutto $\delta$ il rumore copriva il segnale,
e U non si accorgeva delle inversioni.

\voce{Se mancasse (lesione \lesione{U}; la lesione \lesione{R} toglie solo l'arbitraggio)}
La decisione avrebbe sempre lo stesso equilibrio tra veloce e lento. Dopo un'inversione
l'abitudine continuerebbe a pesare come prima, e il recupero sarebbe più lento. In un
mondo stabile non si noterebbe niente: U serve solo quando il mondo cambia.

\voce{Nel cervello}
Il locus coeruleus e la noradrenalina, che segnalano l'incertezza inattesa (Yu e Dayan,
2005), con la corteccia cingolata anteriore che segue la volatilità (Behrens e colleghi,
2007). L'arbitraggio è il passaggio tra sistema orientato allo scopo e abitudini (Daw, Niv,
Dayan, 2005).

\voce{Misure}
Arbitraggio: circa +3 punti nel recupero (50,0\% contro 47,1\% a 60--90 s). L'effetto è
piccolo perché le sorprese sono frequenti anche prima di un'inversione, quindi
un'inversione non produce un'ondata riconoscibile.

% --- soglia ---
\subsection{La soglia sulla sorpresa}

\voce{Struttura}
Non ha neuroni: è un calcolo tra il corpo e chi impara. Ha 2 ingressi ($\delta$ e la sua media), 1 uscita $\hat\delta$ e
due geni: \gene{floor} = 1 e \gene{u\_slow} = 200 s. Limite fisso: $|\hat\delta|\le3$.

\voce{Come funziona, a ogni passo}
\begin{enumerate}
\item Il corpo dà il valore del momento:
\[
r=\text{piacere}\,\bigl(1+0{,}8\,(2\,\text{fame}-1)\bigr)-\text{dolore}-\text{energia spesa},\quad\text{fame}=1-\text{energia},
\]
  tutto misurato in pile buone (una pila buona = 1 in 2 s, cioè 0,025 a passo). A questo si
  sommerebbero i valori innati sui sensi, che nella fase A sono zero.
\item Con \gene{critic\_mix} = 0 la sorpresa è $\delta=r$.
\item Aggiorna la sorpresa solita, $\overline{|\delta|}$, con $\tau=200$ s.
\item Calcola la parte che insegna:
\[
\hat\delta=\operatorname{sign}(\delta)\;\frac{\max\bigl(|\delta|-\overline{|\delta|},0\bigr)}{\overline{|\delta|}},
\qquad |\hat\delta|\le3.
\]
\end{enumerate}
Misurando in una corsa, $|\delta|$ ha mediana 0,0007 e la sorpresa solita è 0,0013. Muoversi quindi resta sotto la
soglia e non insegna niente. Mangiare ($\approx0{,}025$) o il dolore ($0{,}32$) stanno decine di volte sopra, e arrivano tutti
al tetto 3: il dolore pesa quanto il cibo. È il compromesso descritto sotto.

\voce{Perché serve}
Nel mondo c'è sempre un po' di $\delta$: spendere energia è un piccolo dispiacere continuo.
Senza soglia questo rumore cambierebbe tutti i pesi a ogni passo. La scala relativa fa
sì che E, C e il cancello imparino allo stesso modo con premi grandi o piccoli.

\voce{Se mancasse}
Si tornerebbe al rumore che copriva il segnale nelle reti di prima. Attenzione però: anche
qui c'è un compromesso misurato. Il tetto schiaccia le sorprese grandi e rare, come il
dolore di una cattiva. Senza la scala adattata si sceglie meglio (66\%) ma si mangia meno
(2,46 pile al minuto contro 3,54) e alla fine si ottiene meno energia (33,9 contro 40,5).

\voce{Nel cervello}
I neuroni a dopamina hanno un livello di base e adattano la scala alla distribuzione dei
premi attesi (Tobler, Fiorillo, Schultz, 2005).

% --- F ---
\subsection{F --- Motore}

\voce{Struttura}
\begin{struttura}
Unità & 4 neuroni motori: gira a sinistra, gira a destra, avanti, indietro. Attività $a=[\sum_i W_i x_i]_0^1$. \\
Ingressi & 14, calcolati dal programma e dall'oggetto (tabella sotto). \\
Pesi & $4\times14=56$, fissi in vita, geni. Di partenza solo 13 non sono zero. Il neurone ``indietro'' ha tutti i pesi a zero. \\
Stato & \texttt{wander}: la deriva del vagare. \texttt{escape}: da che parte si sta girando per fuggire da un muro. \\
Uscite & $\text{spinta}=a_{\text{avanti}}-a_{\text{indietro}}$; $\text{sterzo}=a_{\text{destra}}-a_{\text{sinistra}}$. \\
\end{struttura}

\begin{center}
\small
\begin{tabularx}{\textwidth}{@{}rYlr@{}}
\toprule
\textbf{N.} & \textbf{Ingresso} & \textbf{Va a} & \textbf{Peso} \\
\midrule
1, 2 & avvicinati: oggetto a sinistra, a destra & gira da quella parte & 3,1 \\
3 & avvicinati: $1-$vicino & avanti & 1,35 \\
4 & avvicinati: vicino & avanti & 0,15 \\
5, 6 & allontanati: oggetto a sinistra, a destra & gira dalla parte opposta & 2,0 \\
7 & allontanati & avanti & 0,7 \\
8 & vaga & avanti & 0,8 \\
9, 10 & vaga: deriva a sinistra, a destra & gira da quella parte & 1,0 \\
11, 12 & riflesso: muro, gira a sinistra, a destra & gira da quella parte & 4,0 \\
13 & muro davanti: $\text{raggio centrale}-0{,}3$ & avanti & $-0{,}7$ \\
14 & costante & nessuno & 0 \\
\bottomrule
\end{tabularx}
\end{center}

\voce{Come funziona, a ogni passo}
\begin{enumerate}
\item La deriva fa un passo a caso: $\text{wander}\leftarrow0{,}9\,\text{wander}+0{,}3\,\mathcal N(0,1)$.
\item Riflesso: se un raggio supera 0,5 c'è un muro; si sceglie la parte più libera e ci si resta finché davanti non è libero.
  Senza questa memoria, negli angoli si oscillava per sempre.
\item Si calcolano i tre interruttori: \emph{avvicinati} se il programma è 0 e c'è un oggetto; \emph{allontanati} se è 1 e c'è un
  oggetto; \emph{vaga} in tutti gli altri casi. ``Avvicinati'' senza oggetto, ``lascialo'' ed ``esplora'' sono quindi lo stesso movimento.
\item Si riempiono i 14 ingressi moltiplicando gli interruttori per la posizione dell'oggetto.
\item Si calcolano i 4 neuroni e, da loro, spinta e sterzo. Il mondo li usa: fino a 120 unità al secondo e 3 radianti al secondo.
\end{enumerate}
\textbf{Il morso non passa da F}: è una decisione di E che va direttamente al corpo. Riesce solo se una pila è a meno di
$10+8+6=24$ unità, entro $25^\circ$ dal muso, e l'essere va a meno del 30\% della velocità massima. Costa 0,1 di energia
al secondo. Mentre consuma una pila la spinta è forzata a zero.

\voce{Perché serve}
Separa \emph{che cosa fare} da \emph{come farlo}. Sapere muoversi non deve essere
reimparato, né può essere rovinato da ciò che si impara. Così anche i programmi di E hanno
un significato stabile: ``avvicinati'' fa sempre avvicinare. Il riflesso dei muri non passa
da E: reagisce più in fretta di una decisione ed evita di restare incastrati.

\voce{Se mancasse}
Il cervello saprebbe scegliere ma il corpo non si muoverebbe. Se invece F imparasse, ogni
inversione rischierebbe di rovinare i movimenti, e nei primi minuti di vita gli esseri non
saprebbero nemmeno girarsi verso una pila, quindi non ci sarebbe nessun premio da cui
imparare.

\voce{Nel cervello}
Il tronco encefalico e il midollo spinale, con i generatori di ritmo e i riflessi (come il
ritiro davanti a un ostacolo). \textbf{Manca il cervelletto}: il cervello non ha nessuna
sezione che impari con maestro a correggere i movimenti.

% ---------------------------------------------------------------------------
\section{Il flusso dell'informazione: dai sensi ai motori}\label{sec:flusso}

Questa sezione segue l'informazione invece delle parti: da dove entra, quante strade fa
per arrivare ai motori e dove torna su sé stessa.

\subsection{Che forma ha l'informazione, tappa per tappa}

\begin{center}
\small
\begin{tabularx}{\textwidth}{@{}lrY@{}}
\toprule
\textbf{Tappa} & \textbf{Numeri} & \textbf{Che cosa contiene} \\
\midrule
Sensi & 31 & occhi, raggi, urti, energia, gusto, dolore \\
Oggetti in vista & fino a 13 & per ognuno: caratteristiche, vicinanza, angolo \\
Mappa di S & 9 posizioni & punteggio dell'oggetto più vicino in ogni direzione \\
Fovea & 10 & un oggetto solo \\
A & 52 & sensi, ciò che è appena salito, fovea \\
B & 16 & la situazione: 3 prototipi accesi \\
D & 10 & l'ultimo oggetto guardato \\
C & 256 & il codice dell'oggetto: in media 13,5 neuroni accesi \\
O & 2 & quanto vale ciò che guardo, e ciò che ho preso \\
Ingressi di E & 29 & B, D, O e una costante \\
Voti di E & 5 & 4 programmi e il morso \\
Scelta & 2 & un programma, e sì o no al morso \\
Ingressi di F & 14 & il programma combinato con dove sta l'oggetto \\
Motori & 4 & gira a sinistra, a destra, avanti, indietro \\
Comandi & 3 & sterzo, spinta, morso \\
\bottomrule
\end{tabularx}
\end{center}

L'informazione si restringe molto: da decine di numeri a una scelta tra pochi programmi.
Solo dopo la scelta si riallarga, quando F le ricombina la posizione dell'oggetto.
L'imbuto stretto è voluto: E sceglie \emph{che cosa} fare, F \emph{come} farlo.

\subsection{Le cinque strade verso i motori}

Non esiste una sola catena dai sensi ai motori. Ci sono cinque strade, dalla più rapida e
rigida alla più lenta e flessibile (figura~\ref{fig:strade}).

\begin{figure}[ht]
\centering
\begin{tikzpicture}[x=1cm,y=1cm]
\tikzset{r1/.style={-{Stealth[length=2mm]},line width=.9pt,black!70},
         r2/.style={-{Stealth[length=2mm]},line width=.9pt,orange!85!black},
         r3/.style={-{Stealth[length=2mm]},line width=.9pt,red!70!black},
         r4/.style={-{Stealth[length=2mm]},line width=.9pt,green!45!black},
         r5/.style={-{Stealth[length=2mm]},line width=.9pt,blue!70!black}}
\node[sec,fill=fix,minimum width=15mm] (OB) at (0.8,6.4) {Oggetti\\in vista};
\node[sec,fill=fix,minimum width=15mm] (SE) at (0.8,3.0) {Sensi};
\node[sec,fill=sens,minimum width=17mm] (S) at (3.6,6.4) {\textbf{S}};
\node[sec,fill=sens,minimum width=17mm] (A) at (3.6,3.0) {\textbf{A}};
\node[sec,fill=val,minimum width=17mm] (O) at (7.0,8.3) {\textbf{O}};
\node[sec,fill=sens,minimum width=17mm] (FO) at (7.0,6.4) {Fovea};
\node[sec,fill=mem,minimum width=17mm] (D) at (7.0,4.6) {\textbf{D}};
\node[sec,fill=sens,minimum width=17mm] (B) at (7.0,3.0) {\textbf{B}};
\node[sec,fill=mem,minimum width=17mm] (C) at (10.0,5.3) {\textbf{C}};
\node[sec,fill=dec,minimum width=20mm] (E) at (12.2,3.0) {\textbf{E}};
\node[sec,fill=fix,minimum width=16mm] (F) at (13.6,6.4) {\textbf{F}};
\node[sec,fill=fix,minimum width=15mm] (M) at (15.6,3.0) {Corpo\\e mondo};

\draw[flow] (OB) -- (S);
\draw[flow] (SE) -- (A);
\draw[flow] (S) -- (FO);
\draw[flow] (S) -- (A);
\draw[r1] (SE.south) -- (0.8,0.9) -- node[lab,above] {1. riflesso: i raggi fanno girare via dal muro} (13.9,0.9) -- (13.9,5.85);
\draw[r2] (3.9,6.95) -- (3.9,7.4) -- node[lab,above] {2. guida: dove sta l'oggetto, a ogni passo} (13.6,7.4) -- (F.north);
\draw[r3] (O.west) -| node[lab,pos=.25,above] {3. valore} (3.3,6.95);
\draw[r4] (FO.east) -- (9.7,6.4) -- (9.7,5.85);
\draw[r4] (D.east) -- (8.5,4.6) -- (8.5,5.1) -- (9.15,5.1);
\draw[r4] (FO) -- (D);
\node[lab,text=green!45!black] at (8.95,6.65) {4. memoria veloce};
\draw[r4] (C.east) -| (12.2,3.55);
\draw[r5] (A) -- (B);
\draw[r5] (B) -- node[lab,above] {5. abitudine} (E);
\draw[r5] (7.6,4.05) -- (7.6,3.8) -- (11.5,3.8) -- (11.5,3.55);
\draw[flow] (13.3,3.0) -- (13.3,5.85);
\draw[flow] (F.east) -| (M.north);
\draw[flow] (E.east) ++(0,-.3) -- ([yshift=-3mm]M.west);
\node[lab] at (14.38,2.5) {morso};
\end{tikzpicture}
\caption{Le cinque strade dai sensi ai motori. Le strade 1 e 2 non passano da E: agiscono a
ogni passo, cioè ogni 0,05 s. La strada 3 agisce su che cosa si guarda. Le strade 4 e 5 votano
in E, che sceglie ogni 0,5 s. Il morso è l'unico comando che esce da E senza passare da F.}
\label{fig:strade}
\end{figure}

\begin{enumerate}
\item \textbf{Il riflesso.} I raggi di distanza vanno direttamente in F: un muro vicino fa
  girare dalla parte più libera. Non passa per nessuna sezione che impara e non si può
  disimparare. Agisce allo stesso passo in cui il muro viene percepito.
\item \textbf{La guida.} S dice a F dove sta l'oggetto attento, a ogni passo. E decide
  ``avvicinati'' una volta ogni 0,5 s, ma lo sterzo si corregge di continuo senza che E
  debba intervenire. Così una decisione lenta produce un movimento preciso.
\item \textbf{Il valore.} O alza il punteggio degli oggetti buoni in S. Questa strada non
  arriva ai motori con un comando: cambia il bersaglio delle strade 2, 4 e 5. È la via
  principale con cui il sapere diventa comportamento. Senza S che legge O la scelta
  torna al 52,5\%, quanto senza O.
\item \textbf{La memoria veloce.} L'identità nella fovea e in D accende un codice in C,
  che vota in E. È la strada che impara in un colpo, e pesa di più quando U è alta.
\item \textbf{L'abitudine.} Tutti i sensi passano da A a B e arrivano in E, insieme a D e O.
  È l'unica strada in cui entrano muri, urti e ``mi guardano'' per decidere, ed è lenta.
  Anche il voto di L si somma qui.
\end{enumerate}

Il critico V non compare: legge da tutte le sezioni, ma la sua uscita finisce solo sullo
schermo. Oggi, nel flusso dell'informazione, è un vicolo cieco.

\subsection{Quanto ci mette}

Dentro un passo il calcolo è una sequenza, senza attese. Perciò anche l'ordine dei calcoli
della sezione~\ref{sec:schema} fissa i ritardi.

\begin{center}
\small
\begin{tabularx}{\textwidth}{@{}Yl@{}}
\toprule
\textbf{Da che cosa a che cosa} & \textbf{Ritardo} \\
\midrule
Muro vicino $\to$ sterzata di riflesso & 1 passo (0,05 s) \\
Oggetto che si sposta $\to$ correzione dello sterzo & 1 passo \\
Oggetto nuovo guardato $\to$ nuova decisione di E & 1 passo: E decide subito \\
Altri cambiamenti del mondo $\to$ nuova decisione di E & fino a 0,5 s \\
Azione $\to$ piacere o dolore in ingresso & 1 passo \\
Piacere o dolore $\to$ apprendimento di C ed E & quasi subito, se la sorpresa supera la soglia \\
Esito di una pila $\to$ apprendimento di O e K & a fine episodio: circa 2 s per una pila buona \\
\bottomrule
\end{tabularx}
\end{center}

\subsection{I cicli}

Sì, ci sono cicli, e molti. \textbf{Ma non ci sono cicli istantanei.} Dentro un passo
l'informazione scorre in una sola direzione: ogni sezione viene calcolata una volta, in
ordine. Un ciclo si chiude sempre perché una sezione legge lo stato che un'altra ha lasciato
al passo precedente, quindi con almeno 0,05 s di ritardo. Non c'è una rete ricorrente che si
assesta, come nel cervello: qualcosa di simile esiste solo nel fatto che lo stato di B e di D
porta con sé il passato.

La figura~\ref{fig:cicli} mostra i cicli principali. I numeri sono quelli delle tabelle.

\begin{figure}[ht]
\centering
\begin{tikzpicture}[x=1cm,y=1cm]
\tikzset{num/.style={circle,draw=black!70,fill=white,inner sep=.6pt,font=\scriptsize\bfseries,minimum size=4.2mm},
         cy/.style={-{Stealth[length=2mm]},line width=.8pt,red!70!black},
         wo/.style={-{Stealth[length=2mm]},line width=.8pt,blue!70!black}}
\node[sec,fill=sens,minimum width=16mm] (S) at (1.8,4.2) {\textbf{S}};
\node[sec,fill=sens,minimum width=16mm] (FO) at (5.0,4.2) {Fovea};
\node[sec,fill=mem,minimum width=16mm] (C) at (8.2,4.2) {\textbf{C}};
\node[sec,fill=dec,minimum width=16mm] (E) at (11.4,4.2) {\textbf{E}};
\node[sec,fill=fix,minimum width=16mm] (F) at (14.4,4.2) {\textbf{F}};
\node[sec,fill=val,minimum width=16mm] (O) at (5.0,7.0) {\textbf{O}};
\node[sec,fill=val,minimum width=16mm] (K) at (8.2,7.0) {\textbf{K}};
\node[sec,fill=fix,minimum width=142mm] (M) at (8.2,0.7) {Corpo e mondo};

\draw[flow] (S) -- (FO);
\draw[flow] (FO) -- (C);
\draw[flow] (C) -- (E);
\draw[flow] (E) -- (F);
% 1: O S fovea O
\draw[cy] (O.west) -| (1.8,4.75);
\draw[cy] (FO) -- (O);
\node[num] at (3.0,7.0) {1};
% 2: E S
\draw[cy] ([xshift=4mm]E.north) -- (11.8,8.2) -- (1.0,8.2) -- (1.0,4.75);
\node[num] at (6.4,8.2) {2};
\draw[cy] ([xshift=-4mm]S.south) -- (1.4,3.1) -- (11.0,3.1) -- ([xshift=-4mm]E.south);
\node[num] at (6.4,3.1) {2};
% 3: E C E
\draw[cy] (E.north) ++(-3mm,0) -- (11.1,5.4) -- (8.2,5.4) -- (C.north);
\node[num] at (9.7,5.4) {3};
% 4: K O
\draw[cy] ([yshift=2mm]K.west) -- ([yshift=2mm]O.east);
\draw[cy] ([yshift=-2mm]O.east) -- ([yshift=-2mm]K.west);
\node[num] at (6.6,7.5) {4};
% mondo
\draw[wo] (F.south) -- (14.4,1.25);
\draw[wo] (1.8,1.25) -- (S.south);
\node[num] at (1.8,2.4) {5};
\draw[wo] ([xshift=4mm]E.south) -- (11.8,1.25);
\draw[wo] (M.west) -- (0.3,0.7) -- (0.3,7.75) -- (4.7,7.75) -- (4.7,7.55);
\node[num] at (0.3,5.5) {6};
\node[lab] at (12.5,2.2) {morso};
\end{tikzpicture}
\caption{I cicli principali. In rosso quelli interni al cervello, in blu quelli che passano dal
mondo. \textbf{1}: ciò che vale viene guardato, e ciò che si guarda impara. \textbf{2}: ``lascialo''
cambia l'oggetto; un oggetto cambiato fa decidere subito. \textbf{3}: E marca C, e C rivota per E.
\textbf{4}: K sceglie il cassetto di O, gli errori di O decidono K. \textbf{5}: il movimento
cambia ciò che si vede. \textbf{6}: il morso produce un esito, che insegna.}
\label{fig:cicli}
\end{figure}

\subsubsection*{Cicli interni al cervello}

\begin{center}
\small
\begin{tabularx}{\textwidth}{@{}cL{4.1cm}Y@{}}
\toprule
& \textbf{Percorso} & \textbf{Che cosa produce} \\
\midrule
1 & O $\to$ S $\to$ fovea $\to$ O & \textbf{Si impara su ciò che si guarda, e si guarda ciò che ha valore.} Accelera
  l'apprendimento ma può chiudersi su sé stesso: ciò che non si guarda non si aggiorna. Il termine
  di incertezza serve a riaprirlo. È anche una causa della perseverazione dopo un'inversione. \\
2 & E $\to$ S $\to$ fovea, D, C $\to$ E & Il programma ``lascialo'' fa ignorare l'oggetto per 3 s. L'oggetto
  cambia, e questo fa decidere E subito invece che al prossimo turno. \\
3 & E $\to$ C $\to$ E & L'inerzia: la scelta marca i neuroni accesi di C, e finché dura la marcatura
  quei neuroni votano per la stessa scelta. Tornando su un oggetto riaffiora la decisione presa. \\
4 & K $\to$ O $\to$ K & K sceglie quale cassetto usare; gli errori dei cassetti decidono K. Avviene una
  volta per episodio. \\
\midrule
--- & Anelli su una sola sezione & Lo stato di B segue l'ingresso con $\tau=0{,}2$ s; D tiene il ricordo per 4 s;
  il canale lento di A tiene la media; E parte avvantaggiata sul programma in corso quando non guarda niente. \\
\bottomrule
\end{tabularx}
\end{center}

\subsubsection*{Cicli che passano dal mondo}

\begin{center}
\small
\begin{tabularx}{\textwidth}{@{}cL{4.1cm}Y@{}}
\toprule
& \textbf{Percorso} & \textbf{Che cosa produce} \\
\midrule
5 & Motori $\to$ posizione $\to$ sensi e oggetti & Il ciclo sensomotorio, chiuso a ogni passo. Avvicinarsi
  cambia ciò che si vede; è per questo che la strada della guida funziona senza calcolare traiettorie. \\
6 & Morso $\to$ piacere, dolore $\to$ O, C, E $\to$ scelte & Il ciclo dell'apprendimento: ogni azione produce
  gli esiti da cui si impara la prossima. \\
--- & Sorpresa $\to$ U $\to$ E, C, S $\to$ azioni $\to$ sorpresa & Il ciclo di controllo: tante sorprese spostano
  il peso sulla memoria veloce; se questa funziona le sorprese calano e tutto torna come prima. \\
--- & Movimento $\to$ energia $\to$ fame e riposo & Muoversi consuma energia. L'energia cambia il valore che O
  dà alle cose e dice quando si è a riposo, quindi quando L impara. \\
--- & Io $\to$ compagni $\to$ io & Gli altri esseri mangiano le pile e sono oggetti che si possono guardare. Con il
  maestro O impara anche dai suoi episodi. \\
--- & Vita $\to$ accoppiamento $\to$ geni & Solo nella fase B, il ciclo più lento: chi vive e si accoppia passa i geni, e
  i geni definiscono tutte le costanti e le connessioni. \\
\bottomrule
\end{tabularx}
\end{center}

\subsection{I cicli che mancano}

Sapere dove il flusso \emph{non} torna indietro è utile quanto il resto, perché il cervello
vero è pieno di connessioni di ritorno.
\begin{itemize}
\item \textbf{Niente verso B.} B riceve solo da A. La sua descrizione non può essere guidata da
  ciò che si sta cercando. Solo la velocità di apprendimento dipende da $|\delta|$. L'anello da D ed E
  verso B è nel progetto ma non è stato fatto.
\item \textbf{Niente copia del movimento.} Il cervello non riceve ciò che F comanda e non può
  prevedere le conseguenze dei propri movimenti: se ne accorge solo dai sensi, un passo dopo.
\item \textbf{L non torna in C.} Il sapere va da C a L e non viceversa: L non può rinfrescare
  ciò che C ha dimenticato.
\item \textbf{V non torna a nessuno}, con i valori attuali.
\item \textbf{E usa poco O.} La connessione c'è, ma i pesi restano quasi a zero.
\end{itemize}

% ---------------------------------------------------------------------------
\section{Riassunto: se togliessimo ogni sezione}

Le ultime due colonne tengono separate la teoria e la verifica. ``Non provata'' significa
che nel codice non esiste questa lesione. Le misure vengono da versioni diverse e non sono
state ripetute sul codice di oggi.

\begin{center}
\small
\begin{tabularx}{\textwidth}{@{}lL{3.0cm}L{2.5cm}YL{2.8cm}@{}}
\toprule
& \textbf{Domanda a cui risponde} & \textbf{Nel cervello} & \textbf{Senza, in teoria} & \textbf{Misurato} \\
\midrule
A & Che cosa arriva? & retina, talamo & niente ingressi & non provata \\
S & Che cosa guardo? & collicolo superiore & si guarda ciò che è vicino: il sapere non guida più lo sguardo & scelta da 58\% a 52,5\% \\
B & In che situazione sono? & neocorteccia & si impara sui canali grezzi, più lentamente & nessun effetto chiaro \\
D & Che cosa ho appena preso? & prefrontale & gli esiti tardivi restano senza colpevole & nessun effetto chiaro \\
C & L'ho già visto di recente? & corpo a fungo, ippocampo & riadattamento lento; spegne anche L & 61\% contro 65\% \\
O & Quanto vale questa cosa? & orbitofrontale, amigdala & nessuna preferenza imparata nello sguardo & 58,2\% contro 52,5\% \\
K & In che regime è il mondo? & orbitofrontale, ippocampo & ogni ritorno si reimpara da capo & 5--7 punti al ritorno \\
L & Che cosa vale sempre? & neocorteccia & si perde ciò che resta vero da una vita lunga & nessun effetto chiaro \\
E & Che cosa faccio? & gangli della base & nessun comportamento & non provata \\
V & Che cosa succederà? & striato ventrale, dopamina & oggi nulla; se insegnasse, si imparerebbe solo sul premio & identico, verificato \\
U & Mi fido di ciò che so? & locus coeruleus & recupero più lento dopo un cambio & +3 punti (arbitraggio) \\
F & Come mi muovo? & tronco, midollo & il corpo non si muove & non provata \\
\bottomrule
\end{tabularx}
\end{center}

% ---------------------------------------------------------------------------
\section{Un esempio dall'inizio alla fine}

Un essere nasce mentre le rosse sono buone. Niente di ciò che segue è scritto nel codice.

\begin{enumerate}
\item \textbf{Vede una pila rossa.} S la sceglie perché è nuova e vicina. La fovea si
  riempie; un oggetto nuovo apre il cancello e D ne tiene una copia. In C si accende il
  gruppo di circa 13 neuroni di ``rossa e grande''.
\item \textbf{E decide.} Non sa ancora niente: sceglie quasi a caso e prova ad avvicinarsi
  e a mordere. Le sinapsi dei neuroni accesi si marcano.
\item \textbf{Mangia.} Il gusto sale: comincia un episodio, e viene fotografato ciò che
  aveva in mente. Il piacere è alto perché ha fame: $r$ è grande e supera la soglia.
  C ed E rinforzano le scelte marcate.
\item \textbf{L'episodio finisce.} O alza $w^+$ di rosso e grande. Dalla volta dopo
  una pila rossa ha più valore e S la guarda prima, anche da lontano.
\item \textbf{Dopo qualche minuto} O preferisce la rossa, C vota per morderla, E comincia
  a farne un'abitudine. Da carico e fermo a ricaricarsi, L copia un po' di C.
\item \textbf{Il mondo si inverte.} Morde una rossa: dolore. Il cassetto di K sbaglia
  la previsione e $\delta$ diventa negativo. U sale: C pesa di più, E e L di meno. O sta dicendo ancora
  ``rossa buona'', quindi si sbaglia ancora qualche volta. È la \emph{perseverazione}
  misurata dopo ogni inversione.
\item \textbf{K si accorge.} Nessun cassetto spiega gli esiti: ne viene reclutato uno
  nuovo, e il vecchio resta intatto. O impara la gialla da zero.
\item \textbf{Il mondo torna com'era.} Il cassetto vecchio prevede bene gli esiti, vince
  la prova e l'essere ritrova in pochi episodi ciò che sapeva.
\end{enumerate}

% ---------------------------------------------------------------------------
\section{Il confronto con il cervello}

\subsection{La divisione del lavoro}

La struttura segue Doya (1999): ogni regione usa un tipo diverso di apprendimento. Il
cervello a sezioni aggiunge le memorie veloci e i valori sugli oggetti.

\begin{center}
\begin{tabularx}{\textwidth}{lYY}
\toprule
\textbf{Tipo di apprendimento} & \textbf{Nel cervello} & \textbf{Qui} \\
\midrule
Senza maestro & corteccia & B, L \\
Per rinforzo & gangli della base, dopamina & E, C, cancello di D \\
Con maestro & cervelletto & manca \\
Per associazione con gli esiti & amigdala, orbitofrontale & O, K \\
Memoria rapida & ippocampo & C \\
Memoria per attività & prefrontale & D \\
\bottomrule
\end{tabularx}
\end{center}

\subsection{I modulatori}

Doya (2002) propone che i modulatori regolino i parametri dell'apprendimento. Qui la
corrispondenza è parziale.

\begin{center}
\begin{tabularx}{\textwidth}{lYY}
\toprule
\textbf{Modulatore} & \textbf{Ruolo in teoria} & \textbf{Qui} \\
\midrule
Dopamina & errore di previsione & $\hat\delta$, ma oggi è il valore del momento, non un errore \\
Noradrenalina & incertezza, esplorazione & U, che però sposta il peso delle decisioni \\
Serotonina & quanto lontano si guarda & orizzonti dei critici, fissi nella vita \\
Acetilcolina & velocità di apprendimento & in parte l'accelerazione di U su C \\
\bottomrule
\end{tabularx}
\end{center}

\subsection{Dove l'analogia non regge}

\begin{itemize}
\item \textbf{Quasi nessun anello all'indietro.} Il cervello è pieno di connessioni dall'alto
  al basso. Qui ci sono solo O verso S ed E verso S; non esiste l'anello da D ed E verso B.
\item \textbf{Regioni e sezioni non sono una a una.} L'orbitofrontale appare sia in O sia in
  K, e C è più insetto che mammifero.
\item \textbf{Il maestro non è biologico.} Gli allievi ``guardano dentro'' il maestro: ne
  leggono ciò che aveva in mente e il suo esito. È un'idea di che cosa si trasferisce, non
  un modello dell'osservazione vera.
\item \textbf{La memoria è copiata, non riprodotta.} Il consolidamento è una media, non il
  ripasso notturno dei ricordi.
\item \textbf{Gli oggetti sono regalati.} Il mondo consegna già separati e descritti gli
  oggetti in vista (sezione~\ref{sec:fovea}): la parte più costosa della visione non esiste.
\item \textbf{Niente spazio.} Non c'è niente come l'ippocampo delle cellule di luogo: gli
  esseri non sanno dove sono, e questo pesa molto sul trovare le pile.
\item \textbf{Neuroni continui.} Nessun impulso; molte ``sezioni'' sono poche formule.
\end{itemize}

% ---------------------------------------------------------------------------
\section{A che punto siamo}\label{sec:stato}

\subsection{I numeri}

Misure del banco di prova: fase A, 200 recinti da 6 esseri, 15 minuti, semi 0 e 1.
La fonte è \texttt{docs/TRAGUARDI\_PILE.md}.

\begin{center}
\begin{tabular}{lrrr}
\toprule
& \textbf{Pile buone} & \textbf{Energia/min} & \textbf{Morti di fame/min} \\
\midrule
Guidato a mano (sa qual è la buona) & 94,7\% & 54,7 & 0,000 \\
Cervello a sezioni, versione attuale & 57,8\% & 42,1 & 0,142 \\
Attore-critico classico & 50,4\% & 30,6 & 0,102 \\
\bottomrule
\end{tabular}
\end{center}

Il cervello a sezioni impara davvero: è l'unico dei due che sceglie sopra il caso. Ma è
lontano dal guidato, e muore di fame più dell'attore-critico.

\subsection{I limiti misurati}

\begin{itemize}
\item \textbf{Il sapere arriva poco alle azioni.} O impara bene, ma E lo legge poco: passa
  quasi tutto dall'attenzione.
\item \textbf{L'esperienza non si accumula.} Chi ha vissuto 4 inversioni sceglie il 57\%,
  chi è appena nato il 61\%: C ed E si portano dietro i regimi passati.
\item \textbf{O si riadatta in minuti, non in secondi.}
\item \textbf{Il morso a vuoto.} Senza niente davanti si morde il 60\% delle volte, e
  mordere costa. Spiega parte delle morti di fame.
\item \textbf{B non sa ciò che conta.} I prototipi seguono ciò che è frequente.
\item \textbf{Molte lesioni non si vedono.} B, D ed L non hanno mai mostrato un effetto
  chiaro. Può voler dire che servono compiti più difficili, oppure che non servono: un
  esito negativo qui è un'informazione, e la specifica dice di cambiare la struttura, non
  il contenuto.
\item \textbf{V è spento di fatto} e, con i valori attuali, $\delta$ non è un errore di
  previsione ma il valore grezzo del momento. L'unica vera previsione del sistema è in O.
\end{itemize}

\subsection{Dove questo documento segue il codice e non la specifica}

\texttt{SPEC\_CERVELLO\_SEZIONI.md} descrive in parte la prima versione. Qui vale il codice:
C ha 256 neuroni con 6 ingressi (non 64 con 4); i programmi sono 4 e valgono per l'oggetto
attento (non 5 legati ai colori); V ha tre critici con orizzonti 2, 8, 30 s (non uno
da 4 s); B legge anche la previsione del prossimo ingresso.

% ---------------------------------------------------------------------------
\section*{Riferimenti}
\addcontentsline{toc}{section}{Riferimenti}
\small
\begin{itemize}[leftmargin=1em,itemsep=1pt]
\item Behrens, Woolrich, Walton, Rushworth (2007). Learning the value of information in an uncertain world. \emph{Nature Neuroscience}.
\item Bisley, Goldberg (2010). Attention, intention, and priority in the parietal lobe. \emph{Annual Review of Neuroscience}.
\item Caron, Ruta, Abbott, Axel (2013). Random convergence of olfactory inputs in the \emph{Drosophila} mushroom body. \emph{Nature}.
\item Daw, Niv, Dayan (2005). Uncertainty-based competition between prefrontal and dorsolateral striatal systems. \emph{Nature Neuroscience}.
\item Doya (1999). What are the computations of the cerebellum, the basal ganglia and the cerebral cortex? \emph{Neural Networks}.
\item Doya (2000). Reinforcement learning in continuous time and space. \emph{Neural Computation}.
\item Doya (2002). Metalearning and neuromodulation. \emph{Neural Networks}.
\item Elfwing, Uchibe, Doya, Christensen (2011). Darwinian embodied evolution of the learning ability for survival. \emph{Adaptive Behavior}.
\item Fecteau, Munoz (2006). Salience, relevance, and firing: a priority map for target selection. \emph{Trends in Cognitive Sciences}.
\item Gershman, Blei, Niv (2010). Context, learning, and extinction. \emph{Psychological Review}.
\item Goldman-Rakic (1995). Cellular basis of working memory. \emph{Neuron}.
\item Heald, Lengyel, Wolpert (2021). Contextual inference underlies the learning of sensorimotor repertoires. \emph{Nature}.
\item Kahneman, Tversky (1979). Prospect theory. \emph{Econometrica}.
\item Kahneman, Treisman, Gibbs (1992). The reviewing of object files. \emph{Cognitive Psychology}.
\item Marr (1971). Simple memory: a theory for archicortex. \emph{Philosophical Transactions of the Royal Society B}.
\item McClelland, McNaughton, O'Reilly (1995). Why there are complementary learning systems in the hippocampus and neocortex. \emph{Psychological Review}.
\item Mongillo, Barak, Tsodyks (2008). Synaptic theory of working memory. \emph{Science}.
\item O'Reilly, Frank (2006). Making working memory work. \emph{Neural Computation}.
\item Pearce, Hall (1980). A model for Pavlovian learning. \emph{Psychological Review}.
\item Posner (1980). Orienting of attention. \emph{Quarterly Journal of Experimental Psychology}.
\item Rao, Ballard (1999). Predictive coding in the visual cortex. \emph{Nature Neuroscience}.
\item Rescorla, Wagner (1972). A theory of Pavlovian conditioning. In \emph{Classical Conditioning II}.
\item Schoenbaum, Roesch, Stalnaker, Takahashi (2009). A new perspective on the role of the orbitofrontal cortex in adaptive behaviour. \emph{Nature Reviews Neuroscience}.
\item Schultz, Dayan, Montague (1997). A neural substrate of prediction and reward. \emph{Science}.
\item Tanaka e colleghi (2004). Prediction of immediate and future rewards differentially recruits cortico-basal ganglia loops. \emph{Nature Neuroscience}.
\item Tobler, Fiorillo, Schultz (2005). Adaptive coding of reward value by dopamine neurons. \emph{Science}.
\item Williams (1992). Simple statistical gradient-following algorithms for connectionist reinforcement learning. \emph{Machine Learning}.
\item Wilson, Takahashi, Schoenbaum, Niv (2014). Orbitofrontal cortex as a cognitive map of task space. \emph{Neuron}.
\item Yu, Dayan (2005). Uncertainty, neuromodulation, and attention. \emph{Neuron}.
\end{itemize}

\end{document}
